\chapter{Buku Pedoman Akademik}

\section{Visi dan Misi Program Studi}

\reviewfrompdf{Bagian ini diadopsi dari lampiran pedoman akademik pada template 2020/2021 dan perlu ditinjau bila terdapat perubahan redaksi resmi program studi.}

\subsection*{Visi Program Studi}
Menjadi program studi yang unggul dalam bidang rekayasa perangkat lunak baik di tingkat nasional maupun internasional.

\subsection*{Misi Program Studi}
\begin{enumerate}
\item Melaksanakan pendidikan vokasi yang inovatif berdasarkan pada sistem pendidikan terapan dengan memanfaatkan kemajuan teknologi.
\item Melaksanakan penelitian terapan berbasis produk dan jasa bidang rekayasa perangkat lunak.
\item Melaksanakan pengabdian masyarakat dengan menggunakan kemajuan rekayasa perangkat lunak untuk meningkatkan kesejahteraan.
\item Mewujudkan kerja sama yang saling menguntungkan dengan berbagai pihak pada bidang rekayasa perangkat lunak.
\end{enumerate}

\subsection*{Tujuan Program Studi}
Menghasilkan lulusan bidang rekayasa perangkat lunak yang berketuhanan, beretika dan bermoral baik, berpengetahuan dan berketrampilan tinggi, siap bekerja dan/atau berwirausaha yang mampu bersaing dalam skala nasional dan global.

Menghasilkan penelitian terapan bidang rekayasa perangkat lunak yang berskala nasional dan internasional.

Menghasilkan pengabdian kepada masyarakat melalui penerapan dan penyebarluasan ilmu pengetahuan dan teknologi secara profesional dalam bidang rekayasa perangkat lunak.

Terwujudnya kerja sama yang saling menguntungkan dengan berbagai pihak pada bidang rekayasa perangkat lunak untuk meningkatkan daya saing.

\section{Ringkasan Struktur Kurikulum 2025}

\subsection*{Semester 1}
\begin{tblr}{width=\textwidth,colspec={Q[c,m,0.6cm]Q[l,m,2.2cm]X[l,m]Q[c,m,1.4cm]Q[c,m,1cm]Q[c,m,1cm]Q[c,m,0.8cm]},hlines,vlines,hspan=minimal,rowsep=1pt,row{1}={bg=headgray,font=\bfseries}}
No & Kode MK & Mata Kuliah & Kelompok & Teori & Praktik & SKS \\
1 & RTI251001 & Pancasila & WN & 2 & 0 & 2 \\
2 & RTI251002 & Konsep Teknologi Informasi & WP & 2 & 0 & 2 \\
3 & RTI251003 & Critical thinking dan problem solving & WP & 2 & 0 & 2 \\
4 & RTI251004 & Matematika Dasar & WP & 2 & 0 & 2 \\
5 & RTI251005 & Bahasa Inggris 1 & WP-PT & 2 & 0 & 2 \\
6 & RTI251006 & Dasar Pemrograman & WP & 2 & 0 & 2 \\
7 & RTI251007 & Praktikum Dasar Pemrograman & WP & 0 & 3 & 3 \\
8 & RTI251008 & Keselamatan dan Kesehatan Kerja & WP-PT & 2 & 0 & 2 \\
9 & RTI251009 & Fisika & WP & 2 & 0 & 2 \\
\SetCell[c=6]{r}\textbf{Total SKS} & & & & & & \textbf{19} \\
\end{tblr}

\subsection*{Semester 2}
\begin{tblr}{width=\textwidth,colspec={Q[c,m,0.6cm]Q[l,m,2.2cm]X[l,m]Q[c,m,1.4cm]Q[c,m,1cm]Q[c,m,1cm]Q[c,m,0.8cm]},hlines,vlines,hspan=minimal,rowsep=1pt,row{1}={bg=headgray,font=\bfseries}}
No & Kode MK & Mata Kuliah & Kelompok & Teori & Praktik & SKS \\
1 & RTI252001 & Agama & WN & 2 & 0 & 2 \\
2 & RTI252002 & Aljabar Linier & WP & 2 & 0 & 2 \\
3 & RTI252003 & Desain Antarmuka & WP & 2 & 0 & 2 \\
4 & RTI252004 & Sistem Operasi & WP & 0 & 2 & 2 \\
5 & RTI252005 & Rekayasa Perangkat Lunak & WP & 2 & 0 & 2 \\
6 & RTI252006 & Basis Data & WP & 2 & 0 & 2 \\
7 & RTI252007 & Praktikum Basis Data & WP & 0 & 2 & 2 \\
8 & RTI252008 & Algoritma dan Struktur Data & WP & 2 & 0 & 2 \\
9 & RTI252009 & Praktikum Algoritma dan Struktur Data & WP & 0 & 3 & 3 \\
\SetCell[c=6]{r}\textbf{Total SKS} & & & & & & \textbf{19} \\
\end{tblr}

\subsection*{Semester 3}
\begin{tblr}{width=\textwidth,colspec={Q[c,m,0.6cm]Q[l,m,2.2cm]X[l,m]Q[c,m,1.4cm]Q[c,m,1cm]Q[c,m,1cm]Q[c,m,0.8cm]},hlines,vlines,hspan=minimal,rowsep=1pt,row{1}={bg=headgray,font=\bfseries}}
No & Kode MK & Mata Kuliah & Kelompok & Teori & Praktik & SKS \\
1 & RTI253001 & Manajemen Proyek & WP & 2 & 0 & 2 \\
2 & RTI253002 & Sistem Informasi Manajemen & WP & 2 & 0 & 2 \\
3 & RTI253003 & Kewarganegaraan & WN & 2 & 0 & 2 \\
4 & RTI253004 & Desain dan Pemrograman Web & WP & 0 & 3 & 3 \\
5 & RTI253005 & Basis Data Lanjut & WP & 0 & 3 & 3 \\
6 & RTI253006 & Metode Numerik & WP & 2 & 0 & 2 \\
7 & RTI253007 & Pemrograman Berbasis Objek & WP & 2 & 0 & 2 \\
8 & RTI253008 & Praktikum Pemrograman Berbasis Objek & WP & 0 & 2 & 2 \\
9 & RTI253009 & Bahasa Inggris 2 & WP-PT & 2 & 0 & 2 \\
\SetCell[c=6]{r}\textbf{Total SKS} & & & & & & \textbf{20} \\
\end{tblr}

\subsection*{Semester 4}
\begin{tblr}{width=\textwidth,colspec={Q[c,m,0.6cm]Q[l,m,2.2cm]X[l,m]Q[c,m,1.4cm]Q[c,m,1cm]Q[c,m,1cm]Q[c,m,0.8cm]},hlines,vlines,hspan=minimal,rowsep=1pt,row{1}={bg=headgray,font=\bfseries}}
No & Kode MK & Mata Kuliah & Kelompok & Teori & Praktik & SKS \\
1 & RTI254001 & Sistem Pendukung Keputusan & WP & 0 & 2 & 2 \\
2 & RTI254002 & Analisis dan Desain Berorientasi Objek & WP & 0 & 2 & 2 \\
3 & RTI254003 & Bahasa Indonesia & WN & 2 & 0 & 2 \\
4 & RTI254004 & Kecerdasan Artifisial & WP & 2 & 0 & 2 \\
5 & RTI254005 & Jaringan Komputer & WP & 2 & 0 & 2 \\
6 & RTI254006 & Praktikum Jaringan Komputer & WP & 0 & 2 & 2 \\
7 & RTI254007 & Pemrograman Web Lanjut & WP & 0 & 3 & 3 \\
8 & RTI254008 & Statistik Komputasi & WP & 2 & 0 & 2 \\
9 & RTI254009 & Proyek Sistem Informasi & WP & 0 & 3 & 3 \\
\SetCell[c=6]{r}\textbf{Total SKS} & & & & & & \textbf{20} \\
\end{tblr}

\subsection*{Semester 5}
\begin{tblr}{width=\textwidth,colspec={Q[c,m,0.6cm]Q[l,m,2.2cm]X[l,m]Q[c,m,1.4cm]Q[c,m,1cm]Q[c,m,1cm]Q[c,m,0.8cm]},hlines,vlines,hspan=minimal,rowsep=1pt,row{1}={bg=headgray,font=\bfseries}}
No & Kode MK & Mata Kuliah & Kelompok & Teori & Praktik & SKS \\
1 & RTI255001 & Kewirausahaan Berbasis Teknologi & WP-PT & 2 & 0 & 2 \\
2 & RTI255002 & Metodologi Penelitian & WP & 2 & 0 & 2 \\
3 & RTI255003 & Pemrograman Mobile & WP & 0 & 3 & 3 \\
4 & RTI255004 & Pembelajaran Mesin & WP & 0 & 3 & 3 \\
5 & RTI255005 & Business Intelligence & WP & 0 & 2 & 2 \\
6 & RTI255006 & Penjaminan Mutu Perangkat Lunak & WP-PT & 2 & 0 & 2 \\
7 & RTI255007 & Pengolahan Citra dan Visi Komputer & WP & 0 & 3 & 3 \\
8 & RTI255008 & Administrasi dan Keamanan Jaringan & WP & 0 & 2 & 2 \\
\SetCell[c=6]{r}\textbf{Total SKS} & & & & & & \textbf{19} \\
\end{tblr}

\subsection*{Semester 6}

Semester 6 terdiri dari mata kuliah wajib dan dua jalur pilihan. Mahasiswa memilih salah satu jalur: Reguler atau Magang.

\subsubsection*{Mata Kuliah Wajib}
\begin{tblr}{width=\textwidth,colspec={Q[c,m,0.6cm]Q[l,m,2.2cm]X[l,m]Q[c,m,1.4cm]Q[c,m,1cm]Q[c,m,1cm]Q[c,m,0.8cm]},hlines,vlines,hspan=minimal,rowsep=1pt,row{1}={bg=headgray,font=\bfseries}}
No & Kode MK & Mata Kuliah & Kelompok & Teori & Praktik & SKS \\
1 & RTI256001 & Komunikasi dan Etika Profesi & WP & 2 & 0 & 2 \\
2 & RTI256002 & Pengembangan Karir & WP & 2 & 0 & 2 \\
3 & RTI256003 & Pemrograman Berbasis Framework & WP & 0 & 2 & 2 \\
\end{tblr}

\subsubsection*{Jalur Reguler (Pilihan Reguler)}
\begin{tblr}{width=\textwidth,colspec={Q[c,m,0.6cm]Q[l,m,2.2cm]X[l,m]Q[c,m,1.4cm]Q[c,m,1cm]Q[c,m,1cm]Q[c,m,0.8cm]},hlines,vlines,hspan=minimal,rowsep=1pt,row{1}={bg=headgray,font=\bfseries}}
No & Kode MK & Mata Kuliah & Kelompok & Teori & Praktik & SKS \\
4 & RTI256104 & Proyek Teknologi Terintegrasi & P & 0 & 6 & 6 \\
5 & RTI256105 & Internet of Things & P & 0 & 3 & 3 \\
6 & RTI256106 & Big Data & P & 0 & 2 & 2 \\
7 & RTI256107 & Cloud Computing & P & 0 & 2 & 2 \\
8 & RTI256108 & Komputasi Hijau & P & 2 & 0 & 2 \\
\SetCell[c=6]{r}\textbf{Total SKS Jalur Reguler (wajib + pilihan)} & & & & & & \textbf{21} \\
\end{tblr}

\subsubsection*{Jalur Magang (Pilihan Magang)}
\begin{tblr}{width=\textwidth,colspec={Q[c,m,0.6cm]Q[l,m,2.2cm]X[l,m]Q[c,m,1.4cm]Q[c,m,1cm]Q[c,m,1cm]Q[c,m,0.8cm]},hlines,vlines,hspan=minimal,rowsep=1pt,row{1}={bg=headgray,font=\bfseries}}
No & Kode MK & Mata Kuliah & Kelompok & Teori & Praktik & SKS \\
4 & RTI256204 & Proyek Inovasi & P & 0 & 6 & 6 \\
5 & RTI256205 & Workshop Teknologi Terapan & P & 0 & 3 & 3 \\
6 & RTI256206 & Rekayasa Sistem & P & 0 & 3 & 3 \\
7 & RTI256207 & Teknologi Terapan & P & 2 & 0 & 2 \\
\SetCell[c=6]{r}\textbf{Total SKS Jalur Magang (wajib + pilihan)} & & & & & & \textbf{20} \\
\end{tblr}

\subsection*{Semester 7}
\begin{tblr}{width=\textwidth,colspec={Q[c,m,0.6cm]Q[l,m,2.2cm]X[l,m]Q[c,m,1.4cm]Q[c,m,1cm]Q[c,m,1cm]Q[c,m,0.8cm]},hlines,vlines,hspan=minimal,rowsep=1pt,row{1}={bg=headgray,font=\bfseries}}
No & Kode MK & Mata Kuliah & Kelompok & Teori & Praktik & SKS \\
1 & RTI257001 & Magang & WP & 0 & 20 & 20 \\
\SetCell[c=6]{r}\textbf{Total SKS} & & & & & & \textbf{20} \\
\end{tblr}

\subsection*{Semester 8}
\begin{tblr}{width=\textwidth,colspec={Q[c,m,0.6cm]Q[l,m,2.2cm]X[l,m]Q[c,m,1.4cm]Q[c,m,1cm]Q[c,m,1cm]Q[c,m,0.8cm]},hlines,vlines,hspan=minimal,rowsep=1pt,row{1}={bg=headgray,font=\bfseries}}
No & Kode MK & Mata Kuliah & Kelompok & Teori & Praktik & SKS \\
1 & RTI258001 & Skripsi & WP & 0 & 8 & 8 \\
2 & RTI258002 & Bahasa Inggris Persiapan Kerja & WP-PT & 2 & 0 & 2 \\
\SetCell[c=6]{r}\textbf{Total SKS} & & & & & & \textbf{10} \\
\end{tblr}

\subsection*{Rekapitulasi Total SKS}
\begin{tblr}{colspec={Q[l,m,5cm]Q[c,m,2cm]},hlines,vlines,hspan=minimal,rowsep=1pt,row{1}={bg=headgray,font=\bfseries}}
Semester & SKS \\
Semester 1 & 19 \\
Semester 2 & 19 \\
Semester 3 & 20 \\
Semester 4 & 20 \\
Semester 5 & 19 \\
Semester 6 (Jalur Reguler) & 21 \\
Semester 6 (Jalur Magang) & 20 \\
Semester 7 & 20 \\
Semester 8 & 10 \\
\textbf{Total (Jalur Reguler)} & \textbf{148} \\
\textbf{Total (Jalur Magang)} & \textbf{147} \\
\end{tblr}

\reviewfrompdf{Ditemukan perbedaan bobot SKS antara RPS mata kuliah dan distribusi pada bagian ini untuk 5 mata kuliah berikut. Nilai pada distribusi di atas dipertahankan karena konsisten dengan total SKS tiap semester, namun perlu dikonfirmasi oleh tim kurikulum mata kuliah mana yang sesuai:
\begin{itemize}
\item RTI253005 Basis Data Lanjut: RPS menyatakan 2 SKS, distribusi ini menyatakan 3 SKS.
\item RTI255008 Administrasi dan Keamanan Jaringan: RPS menyatakan 3 SKS, distribusi ini menyatakan 2 SKS.
\item RTI256205 Workshop Teknologi Terapan: RPS menyatakan 6 SKS, distribusi ini menyatakan 3 SKS.
\item RTI256206 Rekayasa Sistem: RPS menyatakan 6 SKS, distribusi ini menyatakan 3 SKS.
\item RTI256207 Teknologi Terapan: RPS menyatakan 6 SKS, distribusi ini menyatakan 2 SKS.
\end{itemize}}

\section{Deskripsi Mata Kuliah}

\reviewfrompdf{Deskripsi mata kuliah berikut disusun langsung dari RPS Kurikulum 2025 (Lampiran II), bukan ditranskripsi dari dokumen kurikulum 2020, karena RPS 2025 adalah sumber yang lebih akurat dan terkini untuk setiap mata kuliah saat ini.}

\begingroup
\small

\subsection*{Semester 1}

\subsubsection*{Pancasila (RTI251001)}
\begin{longtblr}{width=\textwidth,colspec={Q[l,m,3cm]X[l,m]},rowsep=2pt,hlines,vlines,hspan=minimal}
SKS / Jam & 2 SKS/2 jam \\
Semester & 1 \\
CPL & \begin{itemize}\item \textbf{CPL01}: Mampu menginternalisasi nilai ketakwaan, ketaatan hukum, kedisiplinan, profesionalisme, etika profesi, semangat belajar sepanjang hayat, serta kepedulian terhadap isu sosial dan dampak teknologi.\end{itemize} \\
Deskripsi Singkat & Mata kuliah Pancasila membekali mahasiswa untuk memahami sejarah, landasan, filosofi, dan kedudukan Pancasila sebagai dasar Negara Republik Indonesia serta menginternalisasi nilai-nilainya dalam isu HAM, antikorupsi, dan paradigma pembangunan. \\
Bahan Kajian / Materi Pembelajaran & \begin{enumerate}\item Pendidikan Pancasila dalam tinjauan historis, kultural, yuridis, dan filosofis\item Pancasila dalam konteks sejarah perjuangan bangsa Indonesia\item Pancasila sebagai sistem filsafat\item UUD RI 1945 dan amandemen UUD RI 1945\item Trias Politika dan kelembagaan negara menurut UUD RI 1945\item Pancasila sebagai ideologi nasional dan ideologi lain yang berkembang di dunia\item Pancasila dan HAM serta pelaksanaan HAM dalam UUD RI 1945\item Tindak pidana korupsi\item Pancasila sebagai paradigma pembangunan\end{enumerate} \\
Pustaka Utama & \begin{enumerate}\item Modul Ajar Pendidikan Pancasila.\end{enumerate} \\
Pustaka Pendukung & \begin{enumerate}\item Kementerian Riset, Teknologi, dan Pendidikan Tinggi Republik Indonesia. 2016. Buku Ajar Pendidikan Pancasila untuk Perguruan Tinggi. Jakarta: Direktorat Pembelajaran dan Kemahasiswaan Kemenristekdikti.\end{enumerate} \\
\end{longtblr}

\subsubsection*{Konsep Teknologi Informasi (RTI251002)}
\begin{longtblr}{width=\textwidth,colspec={Q[l,m,3cm]X[l,m]},rowsep=2pt,hlines,vlines,hspan=minimal}
SKS / Jam & 2 SKS/4 jam \\
Semester & 1 \\
CPL & \begin{itemize}\item \textbf{CPL02}: Mampu menerapkan prinsip rekayasa sistem dan perangkat lunak untuk menghasilkan solusi aplikasi multi-platform sesuai kebutuhan.\item \textbf{CPL05}: Mampu menghasilkan solusi teknologi komputasi dan infrastruktur digital yang sesuai dengan kebutuhan organisasi dan pengguna.\item \textbf{CPL09}: Mampu menjelaskan cara kerja sistem komputer dan menerapkan pendekatan komputasi untuk menyelesaikan masalah.\end{itemize} \\
Deskripsi Singkat & Dalam mata kuliah ini akan dibahas tentang konsep teknologi, inovasi teknologi, perkembangan IPTEK, etika rekayasa, perkembangan ICT, sistem komputer, konsep sistem komputer, representasi data, aljabar Boolean, flowchart, jaringan komputer dan internet, aplikasi TI di berbagai bidang, serta sertifikasi bidang TI. \\
Bahan Kajian / Materi Pembelajaran & \begin{enumerate}\item Konsep teknologi dan inovasi teknologi\item Perkembangan IPTEK dan etika rekayasa\item Perkembangan ICT\item Sistem komputer dan konsep sistem komputer\item Representasi data, aljabar Boolean, dan flowchart\item Jaringan komputer dan internet\item Aplikasi TI di berbagai bidang dan sertifikasi bidang TI\end{enumerate} \\
Pustaka Utama & \begin{enumerate}\item William Stallings. 2007. Data and Computer Communications. 8th Edition. Pearson Education, Inc. Pearson Prentice Hall.\end{enumerate} \\
Pustaka Pendukung & \begin{enumerate}\item Davis, W.S. Computers and Information Systems: An Introduction. West Publishing Company.\item Situs web dan artikel KTI.\end{enumerate} \\
\end{longtblr}

\subsubsection*{Critical Thinking dan Problem Solving (RTI251003)}
\begin{longtblr}{width=\textwidth,colspec={Q[l,m,3cm]X[l,m]},rowsep=2pt,hlines,vlines,hspan=minimal}
SKS / Jam & 2 SKS/4 jam \\
Semester & 1 \\
CPL & \begin{itemize}\item \textbf{CPL04}: Mampu menyusun deskripsi saintifik hasil kajian, pengembangan, atau implementasi teknologi informasi dan komunikasi dalam bentuk skripsi, laporan tugas akhir, atau artikel ilmiah.\end{itemize} \\
Deskripsi Singkat & Pada mata kuliah ini mahasiswa akan belajar mengenai konsep berpikir kritis dan bagaimana penerapannya dalam menanggapi informasi dalam kehidupan sehari-hari. Selain itu pada mata kuliah ini akan diajarkan mengenai konsep dan teknik-teknik dalam menyelesaikan kasus/permasalahan baik dalam lingkup ujian/tes, maupun pada masalah-masalah yang dijumpai dalam kehidupan sehari-hari. \\
Bahan Kajian / Materi Pembelajaran & \begin{enumerate}\item Berpikir dan Menalar\item Pondasi Berpikir Kritis\item Kemampuan Dasar Pemecahan Masalah\item Applied Critical Thinking\item Kemampuan Pemecahan Masalah Tingkat Lanjut\item Teknik-teknik Lanjutan untuk Pemecahan Masalah\item Penalaran Kritis\end{enumerate} \\
Pustaka Utama & \begin{enumerate}\item Thinking Skills: Critical Thinking and Problem Solving. Second Edition.\end{enumerate} \\
Pustaka Pendukung & \begin{enumerate}\item Critical Thinking Skills For Dummies.\end{enumerate} \\
\end{longtblr}

\subsubsection*{Matematika Dasar (RTI251004)}
\begin{longtblr}{width=\textwidth,colspec={Q[l,m,3cm]X[l,m]},rowsep=2pt,hlines,vlines,hspan=minimal}
SKS / Jam & 2 SKS/4 jam \\
Semester & 1 \\
CPL & \begin{itemize}\item \textbf{CPL09}: Mampu menjelaskan cara kerja sistem komputer dan menerapkan pendekatan komputasi untuk menyelesaikan masalah.\end{itemize} \\
Deskripsi Singkat & Mata kuliah Matematika Dasar mempelajari matematika diskrit, yaitu cabang aljabar yang mengkaji perhitungan bilangan diskrit (tidak kontinyu) dan cara paling efisien menemukan solusinya; dalam konteks teknologi informasi, bilangan biner mendapatkan perhatian utama. \\
Bahan Kajian / Materi Pembelajaran & \begin{enumerate}\item Proposisi dan logika\item Teori himpunan\item Relasi dan fungsi\item Teori bilangan dan sistem bilangan\item Induksi dan rekursi\item Aljabar Boolean\item Kombinatorial\item Teori graf dan pohon\end{enumerate} \\
Pustaka Utama & \begin{enumerate}\item Yan Watequlis, Cahya Rahmad, Deasy Sandhya Elya. 2017. Matematika Diskrit. Polinema Press.\end{enumerate} \\
Pustaka Pendukung & \begin{enumerate}\item Steven G. Krantz. 2009. Discrete Mathematics Demystified. McGraw-Hill.\item Kenneth H. Rosen. 1999. Discrete Mathematics and Its Application. McGraw-Hill.\item C.L. Liu. 1985. Element of Discrete Mathematics. McGraw-Hill, Inc.\item Munir, Rinaldi. 2012. Matematika Diskrit Ed. Revisi Ke-5. Informatika Bandung.\end{enumerate} \\
\end{longtblr}

\subsubsection*{Bahasa Inggris 1 (RTI251005)}
\begin{longtblr}{width=\textwidth,colspec={Q[l,m,3cm]X[l,m]},rowsep=2pt,hlines,vlines,hspan=minimal}
SKS / Jam & 2 SKS/4 jam \\
Semester & 1 \\
CPL & \begin{itemize}\item \textbf{CPL03}: Mampu mengelola tim dan diri sendiri, berkomunikasi secara lisan dan tulisan, serta mempresentasikan gagasan dan hasil kerja secara efektif dan profesional.\end{itemize} \\
Deskripsi Singkat & Mata kuliah Bahasa Inggris 1 melatih kemampuan dan kecakapan mahasiswa dalam Listening, Speaking, Reading, dan Writing secara terintegrasi dalam konteks Teknik Informatika. Topik-topik materi disesuaikan dengan bidang Teknik Informatika yang dapat diterapkan dalam kehidupan sehari-hari dan dunia kerja. Aktivitas belajar meliputi ceramah, diskusi, role play, presentasi, serta project individu maupun kelompok, dan mendukung penerapan Case Method dan Project Based Learning. \\
Bahan Kajian / Materi Pembelajaran & \begin{enumerate}\item Computer Applications: kinds of applications, Simple Present Tense, imperatives, sequencers, software installation process\item Computer Architecture: types of computers, hardware and its functions, computer ads, adjectives, comparatives, superlatives\item Multimedia: multimedia hardware, toolbox of a graphic program, passive sentences, time clauses\item Computer Network: network hardware, network topology, if-clause\item Websites: features, types, analytics, data visualization, characteristics of a good website\item Careers in IT: IT jobs and responsibilities, modals for skills and qualifications\item IT Support Staff: computer problems and solutions, service report, help desk ticket, email\item Workstation Health and Safety: problems, expressions for giving advice and declaring prohibition\end{enumerate} \\
Pustaka Utama & \begin{enumerate}\item Asri, Atiqah Nurul. 2025. English for Informatics 1: 8th Edition. Modul belum dipublikasikan.\end{enumerate} \\
Pustaka Pendukung & \begin{enumerate}\item Esteras, Santiago Remacha. (2010). Infotech English for Computer Users Workbook. Cambridge: Cambridge University Press.\item Esteras, Santiago Remacha. (2011). Infotech English for Computer Users Teacher's Book. Cambridge: Cambridge University Press.\item Glendinning, Eric H \& McEwan, John. (2012). Basic English for Computing Revised and Updated. Oxford: Oxford University Press.\item Olejniczak, Maja. (2011). English for Information Technology 1 Vocational English Course Book. Essex: Pearson Education Limited.\end{enumerate} \\
\end{longtblr}

\subsubsection*{Dasar Pemrograman (RTI251006)}
\begin{longtblr}{width=\textwidth,colspec={Q[l,m,3cm]X[l,m]},rowsep=2pt,hlines,vlines,hspan=minimal}
SKS / Jam & 2 SKS/4 jam \\
Semester & 1 \\
CPL & \begin{itemize}\item \textbf{CPL07}: Mampu menerapkan sistem komputasi cerdas dan pengolahan data untuk mendukung penyelesaian masalah komputasi.\item \textbf{CPL09}: Mampu menjelaskan cara kerja sistem komputer dan menerapkan pendekatan komputasi untuk menyelesaikan masalah.\end{itemize} \\
Deskripsi Singkat & Dasar Pemrograman memberikan pengetahuan dan pemahaman konsep dasar algoritma dan dasar pemrograman sehingga mahasiswa memiliki dasar untuk menyelesaikan permasalahan-permasalahan logika dengan menggunakan flowchart dan pseudocode. \\
Bahan Kajian / Materi Pembelajaran & \begin{enumerate}\item Konsep algoritma, analisis permasalahan, version control, dan Kanban\item Tipe data, variabel, konstanta, nilai, ekspresi, input-output, sequence, dan operator\item Analisa kasus: pemilihan dan perulangan (sederhana dan bersarang)\item Array 1 dan 2 dimensi, fungsi iteratif dan rekursif\end{enumerate} \\
Pustaka Utama & \begin{enumerate}\item Sebesta, Robert. 2016. Concept of Programming Languages, Edisi Global. Addison Wesley.\item Sestoft, Peter. 2017. Programming Language Concepts. Springer.\item T. Henny Febriana Harumy. 2016. Belajar Dasar Algoritma dan Pemrograman C++. Deepublish.\end{enumerate} \\
Pustaka Pendukung &  Rinaldi Munir. 2015. Algoritma dan Pemrograman. Penerbit Informatika. \\
\end{longtblr}

\subsubsection*{Praktikum Dasar Pemrograman (RTI251007)}
\begin{longtblr}{width=\textwidth,colspec={Q[l,m,3cm]X[l,m]},rowsep=2pt,hlines,vlines,hspan=minimal}
SKS / Jam & 3 SKS/6 jam \\
Semester & 1 \\
CPL & \begin{itemize}\item \textbf{CPL07}: Mampu menerapkan sistem komputasi cerdas dan pengolahan data untuk mendukung penyelesaian masalah komputasi.\item \textbf{CPL09}: Mampu menjelaskan cara kerja sistem komputer dan menerapkan pendekatan komputasi untuk menyelesaikan masalah.\end{itemize} \\
Deskripsi Singkat & Mata kuliah ini membahas pembuatan algoritma untuk menyelesaikan masalah dengan metode yang efisien hingga menghasilkan flowchart dan mentranslasikan teks algoritma ke dalam teks program bahasa pemrograman Java melalui praktik jobsheet terstruktur. \\
Bahan Kajian / Materi Pembelajaran & \begin{enumerate}\item Konsep algoritma, representasi algoritma, translator, dan bahasa pemrograman\item Tipe data, variabel, konstanta, nilai, ekspresi, dan input-output\item Sequence, analisa kasus, pencabangan, dan perulangan\item Array dan fungsi/prosedur\end{enumerate} \\
Pustaka Utama & \begin{enumerate}\item Sebesta, Robert. 2016. Concept of Programming Languages, Edisi Global. Addison Wesley.\item Sestoft, Peter. 2017. Programming Language Concepts. Springer.\item T. Henny Febriana Harumy. 2016. Belajar Dasar Algoritma dan Pemrograman C++. Deepublish.\end{enumerate} \\
Pustaka Pendukung &  Rinaldi Munir. 2015. Algoritma dan Pemrograman. Penerbit Informatika; jobsheet praktikum dan materi LMS. \\
\end{longtblr}

\subsubsection*{Keselamatan dan Kesehatan Kerja (RTI251008)}
\begin{longtblr}{width=\textwidth,colspec={Q[l,m,3cm]X[l,m]},rowsep=2pt,hlines,vlines,hspan=minimal}
SKS / Jam & 2 SKS/4 jam \\
Semester & 1 \\
CPL & \begin{itemize}\item \textbf{CPL01}: Mampu menerapkan nilai ketakwaan, kesadaran hukum, kedisiplinan, profesionalisme, etika profesi, dan pembelajaran sepanjang hayat, serta tanggap terhadap isu sosial dan perkembangan teknologi.\end{itemize} \\
Deskripsi Singkat & Mahasiswa mengetahui perundang-undangan yang menjadi dasar semua kebijakan yang berhubungan dengan kesehatan dan keselamatan kerja, serta mengerti lingkup kesehatan, lingkungan kerja, keselamatan kerja, asuransi kerja, dan organisasi kerja yang menjadi bagian dari sebuah sistem tenaga kerja. \\
Bahan Kajian / Materi Pembelajaran & \begin{enumerate}\item Konsep K3: sejarah, pengertian, dan tujuan K3\item Undang-undang K3: undang-undang yang melandasi K3 dan peraturan pemerintah\item Kesehatan masyarakat: dasar peraturan, pemeriksaan kesehatan sebelum dan setelah kerja\item Lingkungan kerja fisik dan non fisik\item Keselamatan kerja: faktor yang berpengaruh, sumber bahaya, pencegahan kecelakaan kerja\item Alat-alat keselamatan kerja\item Organisasi K3: maksud dan tujuan organisasi K3\item Asuransi: prinsip dasar, jenis-jenis, dan klaim asuransi\item Penerapan protokol kesehatan di pendidikan dan industri era new normal\item Implementasi K3 untuk menciptakan inovasi revolusi industri 4.0\end{enumerate} \\
Pustaka Utama & \begin{enumerate}\item Budi Harijanto. 2012. Modul Ajar K3.\item Rudi Suardi. 2007. Sistem Manajemen Keselamatan dan Kesehatan Kerja. PPM.\item Nuansa Aulia. 2008. Himpunan Peraturan Perundang-undangan Republik Indonesia Keselamatan dan Kesehatan Kerja (K-3) Disertai dengan Peraturan Perundangan yang Terkait.\item Gempur Santosa. 2004. Manajemen Keselamatan dan Kesehatan Kerja. Prestasi Pustaka.\item Occupational Health and Safety Management Systems (OHSAS 18001:2007) -- Requirements.\item Meyti Eka Apriyani. 2022. Keselamatan dan Kesehatan Kerja. Polinema Press.\end{enumerate} \\
Pustaka Pendukung &  Materi LMS dan peraturan perundang-undangan K3 terbaru. \\
\end{longtblr}

\subsubsection*{Fisika (RTI251009)}
\begin{longtblr}{width=\textwidth,colspec={Q[l,m,3cm]X[l,m]},rowsep=2pt,hlines,vlines,hspan=minimal}
SKS / Jam & 2 SKS/2 jam \\
Semester & 1 \\
CPL & \begin{itemize}\item \textbf{CPL09}: Mampu menjelaskan cara kerja sistem komputer dan menerapkan pendekatan komputasi untuk menyelesaikan masalah.\end{itemize} \\
Deskripsi Singkat & Mata kuliah Fisika memberikan pemahaman dasar mengenai konsep-konsep fisika klasik yang relevan dengan bidang teknik informatika. Topik meliputi pengukuran, mekanika, dinamika, energi, momentum, dan kelistrikan, dengan integrasi kasus-kasus dalam sistem komputer, Internet of Things, dan pemrosesan sinyal digital. \\
Bahan Kajian / Materi Pembelajaran & \begin{enumerate}\item Hakikat dan ruang lingkup fisika\item Pengantar ilmu sains dan metode ilmiah\item Pengukuran dan satuan\item Kinematika gerak lurus dan dua dimensi\item Dinamika partikel\item Usaha dan energi\item Impuls dan momentum\item Medan listrik dan gaya Coulomb\item Potensial listrik\item Kapasitansi dan energi dalam medan listrik\item Arus listrik dan analisis rangkaian listrik DC\end{enumerate} \\
Pustaka Utama & \begin{enumerate}\item Serway, R. A. \& Jewett, J. W. 2014. Physics for Scientists and Engineers with Modern Physics (9th Edition). Brooks/Cole -- Cengage Learning.\item Giancoli, D. C. 2005. Physics for Scientists and Engineers with Modern Physics (4th Edition). Pearson Education.\end{enumerate} \\
Pustaka Pendukung &  Purwanto, S. \& Rusyidi, C. 2020. Fisika Dasar: Teori dan Aplikasi dalam Sains dan Teknik. Penerbit UPI Press; materi LMS dan lembar kerja. \\
\end{longtblr}

\subsection*{Semester 2}

\subsubsection*{Agama (RTI252001)}
\begin{longtblr}{width=\textwidth,colspec={Q[l,m,3cm]X[l,m]},rowsep=2pt,hlines,vlines,hspan=minimal}
SKS / Jam & 2 SKS/2 jam \\
Semester & 2 \\
CPL & \begin{itemize}\item \textbf{CPL01}: Menginternalisasi ketakwaan kepada Tuhan Yang Maha Esa, menaati hukum, disiplin, dan menunjukkan profesionalisme melalui etika profesi, pembelajaran sepanjang hayat, serta respons terhadap isu sosial dan teknologi.\end{itemize} \\
Deskripsi Singkat & Mata kuliah wajib umum (MKU) yang mengkaji Pendahuluan, Konsep Manusia dan Agama, Konsep Tauhid, Manusia dalam Islam, Wawasan Ekologi dalam Islam, Aktualisasi Akhlak, IPTEKS dalam Islam, Pandangan Islam tentang Etos Kerja, Konsep Ekonomi Syariah, serta Pembentukan Masyarakat Islam melalui hakikat keluarga dan fungsi masjid. Mahasiswa diajak memahami ajaran Islam berdasarkan Al-Qur'an dan Hadis yang dihubungkan dengan kehidupan nyata era sekarang, sehingga PAI menjadi pegangan dalam pembentukan karakter yang baik secara individu maupun sosial. \\
Bahan Kajian / Materi Pembelajaran & \begin{enumerate}\item Pendahuluan: visi pendidikan agama, deskripsi mata kuliah, pendekatan kuliah, tata tertib, dan sistem penilaian\item Manusia dan Agama: konsep agama dan Islam, agama kebutuhan manusia, dimensi ajaran Islam, metode memahami Islam, misi agama Islam, masa depan agama\item Konsep Tauhid: tauhid kebutuhan manusia, problem syirik, dampak tauhid\item Konsep Manusia: manusia dalam pandangan Islam, fungsi dan tugas hidup manusia, hakikat kehidupan dunia dan akhirat\item Wawasan Ekologi dalam Islam: hakikat alam semesta, arti dan sifat sunnatullah, manfaat alam semesta, Islam dan wawasan lingkungan\item Aktualisasi Akhlak: makna perjuangan Rasul, aktualisasi misi Rasul, aplikasi sifat-sifat Rasul dalam ilmu teknologi, ibadah dan pembentukan akhlak\item IPTEKS: hakikat ilmu dalam Islam, paradigma iptek, sumber ilmu pengetahuan, aplikasi pola dzikir dan pikir\item Etos Kerja: etos kerja dalam Islam, motivasi kerja dalam Islam, aktualisasi jihad dalam pembangunan\item Ekonomi: pengertian, prinsip, dan etika ekonomi syariah, pemberdayaan umat, manajemen zakat, wakaf, infak, dan sedekah\item Masyarakat Islam: fungsi keluarga dalam Islam, pembentukan keluarga sakinah, masjid sebagai pusat peradaban\end{enumerate} \\
Pustaka Utama & \begin{enumerate}\item Al Qur'an Al Karim dan Qur'an Android.\item Fadloli, Sri Nurkudri, Abd. Chalim. 2018. Pendidikan Agama Islam Pada Perguruan Tinggi Umum. Malang: AM Publishing.\item Kemenristekdikti. 2016. Modul Pendidikan Agama Islam Untuk Perguruan Tinggi. Jakarta: Dirjen Belmawa Kemenristekdikti.\item Samiun Jazuli, Ahzami. 2014. Kehidupan Dalam Pandangan Al-Qur'an. Jakarta: Gema Insani.\item Abdullah, M. Amin. 2012. Islamic Studies di Perguruan Tinggi. Yogyakarta: Pustaka Pelajar.\end{enumerate} \\
Pustaka Pendukung & \begin{enumerate}\item Ibnu Katsir, Tafsir al-Qur'an Online.\item Djakfar, H. Muhammad. 2010. Teologi Ekonomi Membumikan Tita Langit di Ranah Bisnis. Malang: UIN Maliki Press.\item Umar, Nasaruddin. 2014. Islam Fungsional Revitalisasi \& Reaktualisasi Nilai Islam. Jakarta: Gramedia.\end{enumerate} \\
\end{longtblr}

\subsubsection*{Aljabar Linier (RTI252002)}
\begin{longtblr}{width=\textwidth,colspec={Q[l,m,3cm]X[l,m]},rowsep=2pt,hlines,vlines,hspan=minimal}
SKS / Jam & 2 SKS/4 jam \\
Semester & 2 \\
CPL & \begin{itemize}\item \textbf{CPL09}: Memiliki pemahaman yang memadai terkait cara kerja sistem komputer dan berbagai teknik/pendekatan berbasis komputasi untuk memecahkan masalah stakeholder.\end{itemize} \\
Deskripsi Singkat & Mata kuliah Aljabar Linier membekali mahasiswa Teknik Informatika dengan konsep dan keterampilan dasar untuk memodelkan serta menyelesaikan masalah komputasi menggunakan pendekatan matematis. Materi mencakup sistem persamaan linier dan metode eliminasi (Gauss dan Gauss-Jordan), operasi matriks, invers dan transpos, determinan (ekspansi kofaktor/Laplace dan Sarrus), aturan Cramer untuk kasus kecil, operasi vektor di 2D/3D (dot product dan cross product), proyeksi ortogonal, serta eigenvalue dan eigenvector. Pembelajaran dirancang berorientasi Outcome-Based Education (OBE) dan student-centered learning melalui latihan terstruktur, pembahasan masalah, serta kasus-kasus sederhana yang relevan dengan bidang informatika sebagai fondasi untuk mata kuliah lanjutan seperti pembelajaran mesin, visi komputer, dan analisis data. \\
Bahan Kajian / Materi Pembelajaran & \begin{enumerate}\item Permodelan Sistem Persamaan Linier (SPL)\item Notasi matriks SPL dan dasar operasi baris elementer\item Metode Gauss dan Gauss-Jordan\item Operasi matriks\item Inversi matriks\item Determinan\item Aturan Cramer\item Vektor\item Eigenvector dan eigenvalue\item Proyeksi ortogonal\end{enumerate} \\
Pustaka Utama & \begin{enumerate}\item Lay, David C., et al. Linear Algebra and Its Applications. 6th ed., Pearson, 2021.\end{enumerate} \\
Pustaka Pendukung & \begin{enumerate}\item Hefferon, Jim. Linear Algebra. Orthogonal Publishing L3c, 2017.\end{enumerate} \\
\end{longtblr}

\subsubsection*{Desain Antarmuka (RTI252003)}
\begin{longtblr}{width=\textwidth,colspec={Q[l,m,3cm]X[l,m]},rowsep=2pt,hlines,vlines,hspan=minimal}
SKS / Jam & 2 SKS/4 jam \\
Semester & 2 \\
CPL & \begin{itemize}\item \textbf{CPL02}: Mampu menerapkan prinsip rekayasa sistem dan perangkat lunak untuk menghasilkan solusi aplikasi multi-platform sesuai kebutuhan.\item \textbf{CPL06}: Mampu merancang, mengimplementasikan, menguji, melakukan deployment, memelihara, serta menjamin mutu perangkat lunak yang menjawab permasalahan dan memenuhi kebutuhan pengguna.\item \textbf{CPL08}: Mampu mengelola proyek teknologi informasi secara profesional melalui perencanaan, koordinasi, komunikasi, dan optimalisasi sumber daya.\end{itemize} \\
Deskripsi Singkat & Mata kuliah Desain Antarmuka membangun kemampuan merancang UI/UX melalui riset pengguna, design thinking, persona, journey map, wireframe, visual design, design system, prototipe interaktif, motion, dan usability testing. \\
Bahan Kajian / Materi Pembelajaran & \begin{enumerate}\item Pengenalan UI/UX dan design thinking\item Persona, journey map, dan kebutuhan pengguna\item Figma, wireframe, visual design, dan design system\item UX writing, prototyping, UX motion, dan usability testing\end{enumerate} \\
Pustaka Utama & \begin{enumerate}\item Buley, L. 2013. The User Experience Team of One. Rosenfeld Media.\item Young, I. 2008. Mental Models. Rosenfeld Media.\item Frost, B. 2016. Atomic Design.\end{enumerate} \\
Pustaka Pendukung &  Materi LMS, dokumentasi resmi perangkat lunak, dan jobsheet praktikum. \\
\end{longtblr}

\subsubsection*{Sistem Operasi (RTI252004)}
\begin{longtblr}{width=\textwidth,colspec={Q[l,m,3cm]X[l,m]},rowsep=2pt,hlines,vlines,hspan=minimal}
SKS / Jam & 2 SKS/4 jam \\
Semester & 2 \\
CPL & \begin{itemize}\item \textbf{CPL02}: Mampu menerapkan prinsip rekayasa sistem dan perangkat lunak untuk menghasilkan solusi aplikasi multi-platform sesuai kebutuhan.\item \textbf{CPL05}: Mampu menghasilkan solusi teknologi komputasi dan infrastruktur digital yang sesuai dengan kebutuhan organisasi dan pengguna.\item \textbf{CPL09}: Mampu menjelaskan cara kerja sistem komputer dan menerapkan pendekatan komputasi untuk menyelesaikan masalah.\end{itemize} \\
Deskripsi Singkat & Mata kuliah Sistem Operasi memberikan pengalaman praktik untuk memahami cara kerja sistem operasi dan melakukan administrasi Linux dari instalasi sampai layanan, keamanan akses, backup, recovery, dan remastering. \\
Bahan Kajian / Materi Pembelajaran & \begin{enumerate}\item Konsep sistem operasi dan instalasi Linux\item Terminal, I/O, proses, Bash, memori, dan system call\item Permission, ACL, user/group, layanan, aplikasi, backup, recovery\item Remastering sistem operasi berbasis Linux\end{enumerate} \\
Pustaka Utama & \begin{enumerate}\item Silberschatz, A., Galvin, P. B., and Gagne, G. 2018. Operating System Concepts, 10th Edition. Wiley.\item Stallings, W. 2017. Operating Systems: Internals and Design Principles, 9th Edition. Pearson.\end{enumerate} \\
Pustaka Pendukung &  Materi LMS, dokumentasi resmi perangkat lunak, dan jobsheet praktikum. \\
\end{longtblr}

\subsubsection*{Rekayasa Perangkat Lunak (RTI252005)}
\begin{longtblr}{width=\textwidth,colspec={Q[l,m,3cm]X[l,m]},rowsep=2pt,hlines,vlines,hspan=minimal}
SKS / Jam & 2 SKS/4 jam \\
Semester & 2 \\
CPL & \begin{itemize}\item \textbf{CPL02}: Mampu menerapkan prinsip rekayasa sistem dan perangkat lunak untuk menghasilkan solusi aplikasi multi-platform sesuai kebutuhan.\item \textbf{CPL06}: Mampu merancang, mengimplementasikan, menguji, melakukan deployment, memelihara, serta menjamin mutu perangkat lunak yang menjawab permasalahan dan memenuhi kebutuhan pengguna.\item \textbf{CPL08}: Mampu mengelola proyek teknologi informasi secara profesional melalui perencanaan, koordinasi, komunikasi, dan optimalisasi sumber daya.\end{itemize} \\
Deskripsi Singkat & Mata kuliah ini membahas proses pengembangan perangkat lunak secara sistematis, mulai dari analisis kebutuhan, perancangan, implementasi, pengujian, hingga pemeliharaan. Materi mencakup model pengembangan seperti Waterfall, Prototype, dan Agile, serta pemodelan dengan DFD, UML, dan flowchart. \\
Bahan Kajian / Materi Pembelajaran & \begin{enumerate}\item Pengantar RPL dan Software Development Life Cycle (SDLC)\item Waterfall, Prototype, Spiral, dan RUP\item Agile, Scrum, dan Kanban\item Pseudocode dan Flowchart\item Data Flow Diagram dan Context Diagram\item UML: Use Case, Use Case Specification, Activity Diagram, dan Sequence Diagram\item Kebutuhan Perangkat Lunak\item Desain Perangkat Lunak\item Implementasi\item Testing\item Evolusi Perangkat Lunak\end{enumerate} \\
Pustaka Utama & \begin{enumerate}\item Ian Sommerville. 2015. Software Engineering, 10th Edition. Pearson.\item William R. King. 2015. Planning for Information Systems. Routledge.\end{enumerate} \\
Pustaka Pendukung & \begin{enumerate}\item Sprague, R.H. and McNurlin, B.C. 2002. Information Systems Management in Practice, 5th edition. Prentice-Hall.\item Ward, J. et al. 1996. Strategic Planning for Information Systems Practice, 2nd edition. Wiley.\end{enumerate} \\
\end{longtblr}

\subsubsection*{Basis Data (RTI252006)}
\begin{longtblr}{width=\textwidth,colspec={Q[l,m,3cm]X[l,m]},rowsep=2pt,hlines,vlines,hspan=minimal}
SKS / Jam & 2 SKS/4 jam \\
Semester & 2 \\
CPL & \begin{itemize}\item \textbf{CPL02}: Menguasai konsep teoritis bidang pengetahuan mengenai berbagai prinsip rekayasa sistem/perangkat lunak dalam merancang solusi aplikasi multi-platform yang relevan dengan kebutuhan stakeholder.\item \textbf{CPL10}: Memiliki kompetensi untuk menganalisis persoalan kompleks untuk mengidentifikasi solusi bidang informatika/ilmu komputer dengan mempertimbangkan wawasan ilmu multidisiplin.\end{itemize} \\
Deskripsi Singkat & Mata kuliah Basis Data membahas konsep dasar basis data relasional, perancangan basis data menggunakan ERD, pemetaan ke model relasional, serta normalisasi. Mahasiswa juga mempelajari implementasi basis data menggunakan MySQL melalui DDL, DML, dan query SELECT, serta penerapannya dalam studi kasus sesuai kebutuhan organisasi atau bisnis. \\
Bahan Kajian / Materi Pembelajaran & \begin{enumerate}\item Pengantar Basis Data\item Entity Relationship Diagram\item Contextual Data Model\item Physical Data Model dan Relational Model\item Normalisasi Basis Data\item Pengantar MySQL dan DDL\item DML MySQL\item Query SELECT\item SELECT Multi Table\item Studi Kasus (identifikasi kebutuhan, ERD, pemetaan, dan implementasi basis data)\end{enumerate} \\
Pustaka Utama & \begin{enumerate}\item Puspitasari, D. and Hani'ah, M. 2019. Cara Mudah Merancang Basis Data Relasional. Polinema Pers.\item Fathansyah. 2018. Basis Data Revisi Ketiga. Bandung: Penerbit Informatika.\end{enumerate} \\
Pustaka Pendukung & \begin{enumerate}\item Silberschatz, A., Korth, H., Sudarshan, S. 2019. Database System Concepts Seventh Edition. McGraw-Hill.\item Dewantara, R., Chirzah, D., Tamsir, N., Rahmayuna, N., Rachman, S., Muhajirin. 2023. Perancangan Basis Data. Klaten: Penerbit Lakeisha.\item Silva, R. 2023. MySQL Crash Course. San Francisco: William Pollock.\end{enumerate} \\
\end{longtblr}

\subsubsection*{Praktikum Basis Data (RTI252007)}
\begin{longtblr}{width=\textwidth,colspec={Q[l,m,3cm]X[l,m]},rowsep=2pt,hlines,vlines,hspan=minimal}
SKS / Jam & 2 SKS/4 jam \\
Semester & 2 \\
CPL & \begin{itemize}\item \textbf{CPL02}: Menguasai konsep teoritis bidang pengetahuan mengenai berbagai prinsip rekayasa sistem/perangkat lunak dalam merancang solusi aplikasi multi-platform yang relevan dengan kebutuhan stakeholder.\item \textbf{CPL07}: Mampu mengembangkan sistem komputasi cerdas, analisis data dengan menerapkan algoritma dan teknik kecerdasan artifisial secara tepat guna.\item \textbf{CPL10}: Memiliki kompetensi untuk menganalisis persoalan kompleks untuk mengidentifikasi solusi bidang informatika/ilmu komputer dengan mempertimbangkan wawasan ilmu multidisiplin.\end{itemize} \\
Deskripsi Singkat & Mata kuliah Praktikum Basis Data berfokus pada praktik perancangan dan implementasi basis data relasional. Kegiatan meliputi penyusunan ERD, pemetaan ke model relasional, normalisasi tabel, serta implementasi menggunakan MySQL melalui perintah DDL, DML, dan query SELECT. Seluruh materi diterapkan dalam studi kasus yang merepresentasikan kebutuhan organisasi atau bisnis. \\
Bahan Kajian / Materi Pembelajaran & \begin{enumerate}\item Pengantar Basis Data (Konsep Data, Informasi, Basis Data Relasional)\item Pemodelan Data dan ER Diagram\item Perancangan Basis Data Menggunakan ER Diagram\item Pemetaan ERD ke Model Relasional\item Normalisasi Basis Data\item Implementasi Basis Data dengan MySQL (DDL)\item Pengelolaan Data dengan SQL (DML)\item Query Data (SQL-DQL dan Multi Table)\item Studi Kasus Terpadu: Perancangan dan Implementasi Basis Data\end{enumerate} \\
Pustaka Utama & \begin{enumerate}\item Puspitasari, D. and Hani'ah, M., 2019, Cara Mudah Merancang Basis Data Relasional, Polinema Pers.\item Fathansyah, 2018, Basis Data Revisi Ketiga, Bandung, Penerbit Informatika.\end{enumerate} \\
Pustaka Pendukung & \begin{enumerate}\item Silberschatz, A., Korth, H., Sudarshan, S., 2019, Database System Concepts Seventh Edition, McGraw-Hill.\item Dewantara, R., et al., 2023, Perancangan Basis Data, Klaten, Penerbit Lakeisha.\item Silva, R., 2023, MySQL Crash Course, San Francisco, William Pollock.\end{enumerate} \\
\end{longtblr}

\subsubsection*{Algoritma dan Struktur Data (RTI252008)}
\begin{longtblr}{width=\textwidth,colspec={Q[l,m,3cm]X[l,m]},rowsep=2pt,hlines,vlines,hspan=minimal}
SKS / Jam & 2 SKS/4 jam \\
Semester & 2 \\
CPL & \begin{itemize}\item \textbf{CPL02}: Menguasai konsep teoritis bidang pengetahuan mengenai berbagai prinsip rekayasa sistem/perangkat lunak dalam merancang solusi aplikasi multi-platform yang relevan dengan kebutuhan stakeholder.\item \textbf{CPL09}: Memiliki pemahaman yang memadai terkait cara kerja sistem komputer dan berbagai teknik/pendekatan berbasis komputasi untuk memecahkan masalah stakeholder.\end{itemize} \\
Deskripsi Singkat & Matakuliah Algoritma dan Struktur Data memberikan pengetahuan dan keterampilan pembuatan algoritma dan struktur data yang meliputi brute force, divide-conquer, stack, queue, linked list, tree, graph, sorting, dan searching. \\
Bahan Kajian / Materi Pembelajaran & \begin{enumerate}\item Searching\item Sorting\item Queue\item Stack\item Linked List\item Tree\item Graph\item Brute Force\item Divide-Conquer\item Depth First Search\item Breadth First Search\end{enumerate} \\
Pustaka Utama & \begin{enumerate}\item Goodrich, M. T., Tamassia, R., and Goldwasser, M. H. 2014. Data Structures \& Algorithms in Java, 6th Edition. Wiley Global Education.\item Ramadhani, C. 2015. Dasar Algoritma dan Struktur Data dengan Bahasa Java. Yogyakarta: Andi Publisher.\item Nugroho, A. 2008. Algoritma dan Struktur Data dalam Bahasa Java. Yogyakarta: Andi Publisher.\item Hariyanto, B. 2007. Struktur Data. Bandung: Informatika.\end{enumerate} \\
Pustaka Pendukung & \begin{enumerate}\item Munir, R. 2015. Algoritma dan Pemrograman. Bandung: Informatika.\end{enumerate} \\
\end{longtblr}

\subsubsection*{Praktikum Algoritma dan Struktur Data (RTI252009)}
\begin{longtblr}{width=\textwidth,colspec={Q[l,m,3cm]X[l,m]},rowsep=2pt,hlines,vlines,hspan=minimal}
SKS / Jam & 3 SKS/6 jam \\
Semester & 2 \\
CPL & \begin{itemize}\item \textbf{CPL07}: Mampu mengembangkan sistem komputasi cerdas, analisis data dengan menerapkan algoritma dan teknik kecerdasan artifisial secara tepat guna.\item \textbf{CPL09}: Memiliki pemahaman yang memadai terkait cara kerja sistem komputer dan berbagai teknik/pendekatan berbasis komputasi untuk memecahkan masalah stakeholder.\end{itemize} \\
Deskripsi Singkat & Matakuliah Praktikum Algoritma dan Struktur Data memberikan pengetahuan dan keterampilan praktis pembuatan algoritma dan struktur data yang meliputi brute force, divide-conquer, stack, queue, linked list, tree, dan graph serta proses sorting dan searching. \\
Bahan Kajian / Materi Pembelajaran & \begin{enumerate}\item Searching\item Sorting\item Queue\item Stack\item Linked List\item Tree\item Graph\item Brute Force\item Divide-Conquer\item Depth First Search\item Breadth First Search\end{enumerate} \\
Pustaka Utama & \begin{enumerate}\item Goodrich, M. T., Tamassia, R., and Goldwasser, M. H. 2014. Data Structures \& Algorithms in Java, 6th Edition. Wiley Global Education.\item Ramadhani, C. 2015. Dasar Algoritma dan Struktur Data dengan Bahasa Java. Yogyakarta: Andi Publisher.\item Nugroho, A. 2008. Algoritma dan Struktur Data dalam Bahasa Java. Yogyakarta: Andi Publisher.\item Hariyanto, B. 2007. Struktur Data. Bandung: Informatika.\end{enumerate} \\
Pustaka Pendukung & \begin{enumerate}\item Buana, I. S., Nata, G. N. M., and Arnawa, I. K. 2018. Struktur Data. Yogyakarta: Andi Publisher.\item Kadir, A. 2015. Teori dan Aplikasi Struktur Data Menggunakan Java. Yogyakarta: Andi Publisher.\end{enumerate} \\
\end{longtblr}

\subsection*{Semester 3}

\subsubsection*{Manajemen Proyek (RTI253001)}
\begin{longtblr}{width=\textwidth,colspec={Q[l,m,3cm]X[l,m]},rowsep=2pt,hlines,vlines,hspan=minimal}
SKS / Jam & 2 SKS/3 jam \\
Semester & 3 \\
CPL & \begin{itemize}\item \textbf{CPL03}: Mampu mengelola tim dan diri sendiri secara profesional serta mengomunikasikan dan mempresentasikan gagasan maupun hasil kerja secara efektif baik lisan maupun tulisan.\item \textbf{CPL04}: Mampu menyusun deskripsi saintifik hasil kajian, pengembangan, atau implementasi teknologi informasi dan komunikasi dalam bentuk skripsi, laporan, atau artikel ilmiah.\item \textbf{CPL08}: Mampu mengelola proyek teknologi informasi secara profesional melalui perencanaan, koordinasi, komunikasi, dan optimalisasi sumber daya.\end{itemize} \\
Deskripsi Singkat & Mata kuliah Manajemen Proyek memberikan pengetahuan dan keterampilan manajemen proyek sistem informasi: siklus manajemen proyek serta pengelolaan ruang lingkup, jadwal, biaya, kualitas (QMS), SDM, risiko, komunikasi, pengadaan, dan pemangku kepentingan proyek, diterapkan pada proyek kelompok sepanjang semester. \\
Bahan Kajian / Materi Pembelajaran & \begin{enumerate}\item Konsep proyek, fase, dan 10 knowledge area manajemen proyek\item Manajemen integrasi, ruang lingkup, waktu, dan biaya proyek\item Manajemen kualitas, SDM, dan pengadaan proyek\item Manajemen komunikasi, risiko, dan pemangku kepentingan proyek\item Agile (Scrum, Kanban, XP), penutupan proyek, dan simulasi proyek perangkat lunak\end{enumerate} \\
Pustaka Utama & \begin{enumerate}\item Chemuturi, M., \& Cagley, T. M. 2010. Mastering Software Project Management: Best Practices. Auerbach Publications.\item Schwalbe, K. 2009. Information Technology Project Management (6th ed.). Course Technology.\item Munir. 2015. Manajemen Proyek Perangkat Lunak. UPI Press.\end{enumerate} \\
Pustaka Pendukung &  Project Management Institute. 2017. A Guide to the Project Management Body of Knowledge (PMBOK Guide) (6th ed.). PMI; materi e-learning lmsslc.polinema.ac.id. \\
\end{longtblr}

\subsubsection*{Sistem Informasi Manajemen (RTI253002)}
\begin{longtblr}{width=\textwidth,colspec={Q[l,m,3cm]X[l,m]},rowsep=2pt,hlines,vlines,hspan=minimal}
SKS / Jam & 2 SKS/4 jam \\
Semester & 3 \\
CPL & \begin{itemize}\item \textbf{CPL02}: Mampu menerapkan prinsip rekayasa sistem dan perangkat lunak untuk menghasilkan solusi aplikasi multi-platform sesuai kebutuhan.\item \textbf{CPL06}: Mampu merancang, mengimplementasikan, menguji, melakukan deployment, memelihara, serta menjamin mutu perangkat lunak yang menjawab permasalahan dan memenuhi kebutuhan pengguna.\item \textbf{CPL08}: Mampu mengelola proyek teknologi informasi secara profesional melalui perencanaan, koordinasi, komunikasi, dan optimalisasi sumber daya.\end{itemize} \\
Deskripsi Singkat & Mata kuliah Sistem Informasi Manajemen membangun kemampuan memahami konsep dasar, struktur, dan penerapan sistem informasi dalam organisasi, mencakup pengambilan keputusan berbasis informasi, SI strategis, dan pengembangan SIM berbasis SDLC. \\
Bahan Kajian / Materi Pembelajaran & \begin{enumerate}\item Konsep dasar sistem dan informasi\item Struktur dan komponen SIM\item Organisasi dan pengambilan keputusan\item SI strategis dan antar organisasi\item Database, TIK, dan pengembangan SIM\end{enumerate} \\
Pustaka Utama & \begin{enumerate}\item O'Brien, J.A. \& Marakas, G.M. 2011. Management Information Systems. McGraw-Hill.\item Laudon, K.C. \& Laudon, J.P. 2020. Management Information Systems. Pearson.\item McLeod, R. 2007. Sistem Informasi Manajemen. PT Indeks.\end{enumerate} \\
Pustaka Pendukung &  Jurnal ilmiah, materi LMS, dan studi kasus aktual organisasi. \\
\end{longtblr}

\subsubsection*{Kewarganegaraan (RTI253003)}
\begin{longtblr}{width=\textwidth,colspec={Q[l,m,3cm]X[l,m]},rowsep=2pt,hlines,vlines,hspan=minimal}
SKS / Jam & 2 SKS/3 jam \\
Semester & 3 \\
CPL & \begin{itemize}\item \textbf{CPL01}: Bertakwa kepada Tuhan Yang Maha Esa dan mampu menunjukkan sikap religius, menjunjung tinggi nilai kemanusiaan, hukum, disiplin, serta etika profesi; berpartisipasi dalam lifelong learning; dan peka terhadap isu sosial-teknologi.\end{itemize} \\
Deskripsi Singkat & Mata kuliah Kewarganegaraan membangun kesadaran kebangsaan dan cinta tanah air melalui kajian identitas nasional, konstitusi, demokrasi, hak asasi manusia, wawasan nusantara, ketahanan nasional, dan integrasi nasional. Mahasiswa mengembangkan sikap demokratis, disiplin, dan partisipatif berdasarkan sistem nilai Pancasila. \\
Bahan Kajian / Materi Pembelajaran & \begin{enumerate}\item Identitas Nasional\item Negara dan Konstitusi\item Hubungan Negara dan Warga Negara\item Negara Hukum\item Demokrasi\item Hak Asasi Manusia\item Wawasan Nusantara\item Ketahanan Nasional\item Integrasi Nasional\end{enumerate} \\
Pustaka Utama & \begin{enumerate}\item Widaningsih. 2023. Pendidikan Kewarganegaraan. Yogyakarta: Pustaka Egaliter.\end{enumerate} \\
Pustaka Pendukung & \begin{enumerate}\item UPT MKU-Politeknik Negeri Malang. 2016. Pendidikan Kewarganegaraan. Malang: Aditya Media Publishing.\item Azra, A. 2003. Demokrasi Hak Asasi Manusia dan Masyarakat Madani. Jakarta: Prenata Media.\item Russel, B. 2004. Sejarah Filsafat Barat dan Kaitannya dengan Kondisi Sosio-Politik dari Zaman Kuno hingga Sekarang. Yogyakarta: Pustaka Pelajar.\end{enumerate} \\
\end{longtblr}

\subsubsection*{Desain dan Pemrograman Web (RTI253004)}
\begin{longtblr}{width=\textwidth,colspec={Q[l,m,3cm]X[l,m]},rowsep=2pt,hlines,vlines,hspan=minimal}
SKS / Jam & 3 SKS/6 jam \\
Semester & 3 \\
CPL & \begin{itemize}\item \textbf{CPL02}: Mampu menerapkan prinsip rekayasa sistem dan perangkat lunak untuk menghasilkan solusi aplikasi multi-platform sesuai kebutuhan.\item \textbf{CPL06}: Mampu merancang, mengimplementasikan, menguji, melakukan deployment, memelihara, serta menjamin mutu perangkat lunak yang menjawab permasalahan dan memenuhi kebutuhan pengguna.\item \textbf{CPL07}: Mampu menerapkan sistem komputasi cerdas dan pengolahan data untuk mendukung penyelesaian masalah komputasi.\end{itemize} \\
Deskripsi Singkat & Mata kuliah ini membahas dasar desain dan pemrograman web melalui HTML, CSS, JavaScript dasar, PHP, form processing, session/cookie, MySQL, integrasi database, dan hosting. \\
Bahan Kajian / Materi Pembelajaran & \begin{enumerate}\item HTML, CSS, JavaScript dasar, dan konsep client-server\item PHP, form processing, upload, cookie, dan session\item MySQL dan pemrograman database dengan PHP\item Hosting dan proyek aplikasi web sederhana\end{enumerate} \\
Pustaka Utama & \begin{enumerate}\item Beaird, J. The Principles of Beautiful Web Design.\item Duckett, J. HTML and CSS: Design and Build Websites.\item Pratama, A. PHP Uncover.\end{enumerate} \\
Pustaka Pendukung &  Materi LMS, dokumentasi resmi perangkat lunak, dan jobsheet praktikum. \\
\end{longtblr}

\subsubsection*{Basis Data Lanjut (RTI253005)}
\begin{longtblr}{width=\textwidth,colspec={Q[l,m,3cm]X[l,m]},rowsep=2pt,hlines,vlines,hspan=minimal}
SKS / Jam & 2 SKS/4 jam \\
Semester & 3 \\
CPL & \begin{itemize}\item \textbf{CPL02}: Mampu menerapkan prinsip rekayasa sistem dan perangkat lunak untuk menghasilkan solusi aplikasi multi-platform sesuai kebutuhan.\item \textbf{CPL07}: Mampu membangun sistem komputasi cerdas dan menganalisis data berbasis kecerdasan buatan untuk mendukung pengambilan keputusan dan inovasi.\item \textbf{CPL10}: Memiliki kompetensi untuk menganalisis persoalan kompleks untuk mengidentifikasi solusi bidang informatika/ilmu komputer dengan mempertimbangkan wawasan ilmu multidisiplin.\end{itemize} \\
Deskripsi Singkat & Mata kuliah Basis Data Lanjut membangun kompetensi mahasiswa dalam mengelola dan mengoptimalkan basis data menggunakan PostgreSQL. Topik meliputi query lanjutan (sub-query, JOIN, CTE), optimasi dengan indeks, pembuatan function, view, stored procedure, manajemen transaksi dan concurrency, full-text search, JSONB, backup-restore, migrasi dari MySQL, serta integrasi basis data dengan aplikasi dalam studi kasus nyata. \\
Bahan Kajian / Materi Pembelajaran & \begin{enumerate}\item Pengenalan PostgreSQL dan query dasar\item Query lanjutan: sub-query, JOIN, grouping, dan CTE\item Indeks dan optimasi query\item Function, view, materialized view, dan stored procedure\item Transaksi, concurrency, full-text search, dan JSONB\item Backup-restore dan migrasi basis data\item Integrasi basis data dengan aplikasi dan studi kasus\end{enumerate} \\
Pustaka Utama & \begin{enumerate}\item Ramakrishnan, R. \& Gehrke, J. 2003. \textit{Database Management Systems}, 3rd ed. McGraw-Hill.\item PostgreSQL Global Development Group. 2024. \textit{PostgreSQL 16 Documentation}. https://www.postgresql.org/docs/.\item Elmasri, R. \& Navathe, S. B. 2016. \textit{Fundamentals of Database Systems}, 7th ed. Pearson.\end{enumerate} \\
Pustaka Pendukung & \begin{enumerate}\item Silberschatz, A., Korth, H., Sudarshan, S. 2019. \textit{Database System Concepts}, 7th ed. McGraw-Hill.\item Momjian, B. 2001. \textit{PostgreSQL: Introduction and Concepts}. Addison-Wesley.\item Materi LMS, dokumentasi resmi PostgreSQL, dan jobsheet praktikum.\end{enumerate} \\
\end{longtblr}

\subsubsection*{Metode Numerik (RTI253006)}
\begin{longtblr}{width=\textwidth,colspec={Q[l,m,3cm]X[l,m]},rowsep=2pt,hlines,vlines,hspan=minimal}
SKS / Jam & 2 SKS/4 jam \\
Semester & 3 \\
CPL & \begin{itemize}\item \textbf{CPL07}: Mampu membangun sistem komputasi cerdas dan menganalisis data berbasis kecerdasan buatan untuk mendukung pengambilan keputusan dan inovasi.\item \textbf{CPL09}: Mampu memahami cara kerja sistem komputer dan menggunakan pendekatan komputasi untuk penyelesaian masalah secara efektif.\item \textbf{CPL10}: Memiliki kompetensi untuk menganalisis persoalan kompleks untuk mengidentifikasi solusi bidang informatika/ilmu komputer dengan mempertimbangkan wawasan ilmu multidisiplin.\end{itemize} \\
Deskripsi Singkat & Mata kuliah Metode Numerik membekali mahasiswa dengan kemampuan menyelesaikan masalah matematika teknik secara numerik menggunakan komputer. Topik meliputi konsep galat, penyelesaian sistem persamaan linear (Gauss, Gauss-Jordan, Gauss-Seidel), persamaan nonlinear (metode tertutup dan terbuka), diferensiasi numerik, integrasi numerik (Riemann, trapezoidal, Simpson), interpolasi, dan regresi. Implementasi dilakukan menggunakan perangkat lunak komputasi. \\
Bahan Kajian / Materi Pembelajaran & \begin{enumerate}\item Pengantar metode numerik dan galat\item Sistem persamaan linear: Gauss, Gauss-Jordan, Gauss-Seidel\item Persamaan nonlinear: metode tertutup dan terbuka\item Diferensiasi numerik\item Integrasi numerik: Riemann, trapezoidal, Simpson\item Interpolasi dan regresi\end{enumerate} \\
Pustaka Utama & \begin{enumerate}\item Rahmad, E.C., Ikawati, D.S.E., \& Syaifudin, Y.W. 2018. \textit{Metode Numerik}. Polinema Press. ISBN 9786026695888.\item Burden, R.L. \& Faires, J.D. 2011. \textit{Numerical Analysis}, 9th ed. Brooks/Cole Cengage Learning.\item Chapra, S.C. \& Canale, R.P. 2015. \textit{Numerical Methods for Engineers}, 7th ed. McGraw-Hill.\end{enumerate} \\
Pustaka Pendukung & \begin{enumerate}\item Heath, M.T. 2002. \textit{Scientific Computing: An Introductory Survey}, 2nd ed. McGraw-Hill.\item Kreyszig, E. 2011. \textit{Advanced Engineering Mathematics}, 10th ed. Wiley.\item Materi LMS, spreadsheet MS Excel, dan lembar kerja praktikum.\end{enumerate} \\
\end{longtblr}

\subsubsection*{Pemrograman Berbasis Objek (RTI253007)}
\begin{longtblr}{width=\textwidth,colspec={Q[l,m,3cm]X[l,m]},rowsep=2pt,hlines,vlines,hspan=minimal}
SKS / Jam & 2 SKS/4 jam \\
Semester & 3 \\
CPL & \begin{itemize}\item \textbf{CPL02}: Mampu menerapkan prinsip rekayasa sistem dan perangkat lunak untuk menghasilkan solusi aplikasi multi-platform sesuai kebutuhan.\item \textbf{CPL07}: Mampu mengembangkan sistem komputasi cerdas dan melakukan analisis data berbasis AI untuk pengambilan keputusan.\item \textbf{CPL09}: Mampu memahami cara kerja sistem komputer dan menerapkan pendekatan komputasi untuk problem solving.\end{itemize} \\
Deskripsi Singkat & Mata kuliah Pemrograman Berbasis Objek membangun kemampuan merancang dan mengimplementasikan program berorientasi objek menggunakan Java, mencakup enkapsulasi, inheritance, polymorfisme, GUI, dan koneksi database. \\
Bahan Kajian / Materi Pembelajaran & \begin{enumerate}\item Konsep dasar OOP: kelas, objek, enkapsulasi\item Relasi kelas, inheritance, dan overriding\item Interface, kelas abstrak, dan polymorfisme\item GUI dan Java API\item Java GUI dengan database\end{enumerate} \\
Pustaka Utama & \begin{enumerate}\item Deitel, P. \& Deitel, H. 2017. Java How to Program. Pearson.\item Bloch, J. 2018. Effective Java. Addison-Wesley.\item Eckel, B. 2006. Thinking in Java. Prentice Hall.\end{enumerate} \\
Pustaka Pendukung &  Materi LMS, dokumentasi resmi Java API, dan jobsheet praktikum. \\
\end{longtblr}

\subsubsection*{Praktikum Pemrograman Berbasis Objek (RTI253008)}
\begin{longtblr}{width=\textwidth,colspec={Q[l,m,3cm]X[l,m]},rowsep=2pt,hlines,vlines,hspan=minimal}
SKS / Jam & 2 SKS/4 jam \\
Semester & 3 \\
CPL & \begin{itemize}\item \textbf{CPL02}: Mampu menerapkan prinsip rekayasa sistem dan perangkat lunak untuk menghasilkan solusi aplikasi multi-platform sesuai kebutuhan.\item \textbf{CPL07}: Mampu mengembangkan sistem komputasi cerdas dan melakukan analisis data berbasis kecerdasan buatan untuk pengambilan keputusan.\item \textbf{CPL09}: Mampu menjelaskan cara kerja sistem komputer dan menerapkan pendekatan komputasi untuk menyelesaikan masalah.\end{itemize} \\
Deskripsi Singkat & Praktikum Pemrograman Berbasis Objek memberikan pengalaman hands-on mengimplementasikan konsep OOP menggunakan Java, mencakup enkapsulasi, inheritance, polymorfisme, GUI Swing, dan koneksi database JDBC dalam proyek akhir terintegrasi. \\
Bahan Kajian / Materi Pembelajaran & \begin{enumerate}\item Kelas, objek, enkapsulasi, konstruktor, dan relasi kelas\item Inheritance, overriding, overloading, kelas abstrak, dan interface\item Polymorfisme, exception handling, collections, dan generics\item GUI Swing: komponen dasar dan event handling\item Koneksi database JDBC dan aplikasi terintegrasi\end{enumerate} \\
Pustaka Utama & \begin{enumerate}\item Deitel, P. \& Deitel, H. 2017. Java How to Program. Pearson.\item Bloch, J. 2018. Effective Java. Addison-Wesley.\item Dokumentasi Oracle Java SE dan JDBC.\end{enumerate} \\
Pustaka Pendukung &  Materi LMS, dokumentasi resmi Java API, dan jobsheet praktikum. \\
\end{longtblr}

\subsubsection*{Bahasa Inggris 2 (RTI253009)}
\begin{longtblr}{width=\textwidth,colspec={Q[l,m,3cm]X[l,m]},rowsep=2pt,hlines,vlines,hspan=minimal}
SKS / Jam & 2 SKS/4 jam \\
Semester & 3 \\
CPL & \begin{itemize}\item \textbf{CPL01}: Mampu menerapkan nilai ketakwaan, kesadaran hukum, kedisiplinan, profesionalisme, etika profesi, dan pembelajaran sepanjang hayat, serta tanggap terhadap isu sosial dan perkembangan teknologi.\item \textbf{CPL03}: Mampu mengelola tim dan diri sendiri secara profesional serta mengomunikasikan dan mempresentasikan gagasan maupun hasil kerja secara efektif baik lisan maupun tulisan.\end{itemize} \\
Deskripsi Singkat & Bahasa Inggris 2 mengembangkan kemampuan membaca teks teknis, menulis laporan akademik, dan berkomunikasi lisan dalam bahasa Inggris pada konteks profesional dan akademik bidang TI. \\
Bahan Kajian / Materi Pembelajaran & \begin{enumerate}\item Reading: kosakata teknis IT, pemahaman teks teknis, dan membaca literatur\item Writing: paragraf akademik, laporan teknis, dan CV serta email profesional\item Listening: ceramah dan presentasi teknis\item Speaking: diskusi, presentasi teknis, dan wawancara simulasi\end{enumerate} \\
Pustaka Utama & \begin{enumerate}\item Murphy, R. 2019. English Grammar in Use. Cambridge University Press.\item Glendinning, E. 2008. Professional English in Use IT. Cambridge University Press.\item Swan, M. 2016. Practical English Usage. Oxford University Press.\end{enumerate} \\
Pustaka Pendukung &  Materi LMS, artikel teknis IT berbahasa Inggris, dan sumber daring terpercaya. \\
\end{longtblr}

\subsection*{Semester 4}

\subsubsection*{Sistem Pendukung Keputusan (RTI254001)}
\begin{longtblr}{width=\textwidth,colspec={Q[l,m,3cm]X[l,m]},rowsep=2pt,hlines,vlines,hspan=minimal}
SKS / Jam & 2 SKS/4 jam \\
Semester & 4 \\
CPL & \begin{itemize}\item \textbf{CPL07}: Mampu mengembangkan sistem komputasi cerdas dan melakukan analisis data berbasis kecerdasan buatan untuk pengambilan keputusan.\item \textbf{CPL10}: Mampu menganalisis persoalan kompleks dengan wawasan multidisiplin untuk menghasilkan solusi di bidang informatika dan ilmu komputer.\end{itemize} \\
Deskripsi Singkat & Mata kuliah Sistem Pendukung Keputusan membangun kemampuan mahasiswa dalam merancang dan mengimplementasikan sistem pendukung keputusan berbasis metode Multi-Criteria Decision Making (MCDM) meliputi WSM, WPM, AHP, ROC, EDAS, MARCOS, MEREC, metode fuzzy, dan GDSS, serta mengevaluasi kinerja sistem DSS yang dibangun. \\
Bahan Kajian / Materi Pembelajaran & \begin{enumerate}\item Konsep SPK, teori keputusan, dan MCDM dasar (WSM, WPM)\item Pembobotan subjektif: AHP (benefit dan cost)\item Metode perankingan lanjutan: ROC, EDAS, MARCOS, MEREC\item GDSS dan metode fuzzy untuk ketidakpastian\item Perancangan, implementasi, dan evaluasi sistem DSS\end{enumerate} \\
Pustaka Utama & \begin{enumerate}\item Turban, E., Sharda, R., Delen, D. 2011. Decision Support and Business Intelligence Systems. 9th ed. Prentice Hall.\item Saaty, T.L. 1980. The Analytic Hierarchy Process. McGraw-Hill.\item Kusumadewi, S., Hartati, S., Harjoko, A., Wardoyo, R. 2006. Fuzzy Multi-Attribute Decision Making (Fuzzy MADM). Graha Ilmu.\end{enumerate} \\
Pustaka Pendukung &  Materi LMS, modul praktikum Excel/Python, dan jurnal terkait metode MCDM dan DSS. \\
\end{longtblr}

\subsubsection*{Analisis dan Desain Berorientasi Objek (RTI254002)}
\begin{longtblr}{width=\textwidth,colspec={Q[l,m,3cm]X[l,m]},rowsep=2pt,hlines,vlines,hspan=minimal}
SKS / Jam & 2 SKS/4 jam \\
Semester & 4 \\
CPL & \begin{itemize}\item \textbf{CPL02}: Mampu menerapkan prinsip rekayasa sistem dan perangkat lunak untuk menghasilkan solusi aplikasi multi-platform sesuai kebutuhan.\item \textbf{CPL06}: Mampu merancang, mengimplementasikan, menguji, melakukan deployment, memelihara, serta menjamin mutu perangkat lunak yang menjawab permasalahan dan memenuhi kebutuhan pengguna.\item \textbf{CPL10}: Mampu menganalisis persoalan kompleks dengan wawasan multidisiplin untuk menghasilkan solusi di bidang informatika dan ilmu komputer.\end{itemize} \\
Deskripsi Singkat & Mata kuliah ini membangun kemampuan menganalisis kebutuhan sistem dan merancang solusi perangkat lunak berorientasi objek menggunakan berbagai diagram UML, mengikuti Software Development Life Cycle (SDLC), untuk menghasilkan Spesifikasi Kebutuhan Perangkat Lunak (SKPL) yang implementable. \\
Bahan Kajian / Materi Pembelajaran & \begin{enumerate}\item Paradigma OO dan pemodelan kebutuhan: domain model, use case, use case description\item Robustness analysis dan activity diagram\item Pemutakhiran domain model, sequence diagram, dan class diagram\item Advanced class diagram, state machine diagram, dan deployment diagram\item Penyusunan dan review Spesifikasi Kebutuhan Perangkat Lunak (SKPL)\end{enumerate} \\
Pustaka Utama & \begin{enumerate}\item Hamilton, K., \& Miles, R. 2006. Learning UML 2.0. O'Reilly.\end{enumerate} \\
Pustaka Pendukung & Hunt, J. 2000. The Unified Process for Practitioners: Object-Oriented Design, the UML and Java. Springer. Lee, M., Kim, H., Kim, J., Lee, J., \& Gum, D. 2005. StarUML 5.0 User Guide. Materi LMS dan jobsheet analisis desain. \\
\end{longtblr}

\subsubsection*{Bahasa Indonesia (RTI254003)}
\begin{longtblr}{width=\textwidth,colspec={Q[l,m,3cm]X[l,m]},rowsep=2pt,hlines,vlines,hspan=minimal}
SKS / Jam & 2 SKS/4 jam \\
Semester & 4 \\
CPL & \begin{itemize}\item \textbf{CPL03}: Mampu mengelola tim dan diri sendiri secara profesional serta mengomunikasikan dan mempresentasikan gagasan maupun hasil kerja secara efektif baik lisan maupun tulisan.\item \textbf{CPL04}: Mampu menyusun deskripsi saintifik hasil kajian, pengembangan, atau implementasi teknologi informasi dan komunikasi dalam bentuk skripsi, laporan, atau artikel ilmiah.\end{itemize} \\
Deskripsi Singkat & Bahasa Indonesia mengembangkan kemampuan berbahasa Indonesia baku untuk keperluan akademik dan profesional, meliputi penulisan laporan teknis, karya tulis ilmiah, dan presentasi di bidang TI. \\
Bahan Kajian / Materi Pembelajaran & \begin{enumerate}\item Ragam bahasa, ejaan, diksi, dan kalimat efektif\item Paragraf, laporan teknis, proposal, dan presentasi ilmiah\item Karya tulis ilmiah: struktur, rumusan masalah, metodologi\item Sitasi, daftar pustaka, abstrak, dan artikel ilmiah\end{enumerate} \\
Pustaka Utama & \begin{enumerate}\item Badan Pengembangan Bahasa. 2022. Pedoman Umum Ejaan Bahasa Indonesia (PUEBI Edisi V).\item Suparno \& Yunus, M. 2009. Keterampilan Dasar Menulis. UT.\item Dalman. 2018. Menulis Karya Ilmiah. Rajawali Press.\end{enumerate} \\
Pustaka Pendukung &  Materi LMS, dokumentasi resmi, dan jobsheet praktikum. \\
\end{longtblr}

\subsubsection*{Kecerdasan Artifisial (RTI254004)}
\begin{longtblr}{width=\textwidth,colspec={Q[l,m,3cm]X[l,m]},rowsep=2pt,hlines,vlines,hspan=minimal}
SKS / Jam & 2 SKS/4 jam \\
Semester & 4 \\
CPL & \begin{itemize}\item \textbf{CPL07}: Mampu mengembangkan sistem komputasi cerdas dan melakukan analisis data berbasis kecerdasan buatan untuk pengambilan keputusan.\item \textbf{CPL09}: Mampu menjelaskan cara kerja sistem komputer dan menerapkan pendekatan komputasi untuk menyelesaikan masalah.\item \textbf{CPL10}: Mampu menganalisis persoalan kompleks dengan wawasan multidisiplin untuk menghasilkan solusi di bidang informatika dan ilmu komputer.\end{itemize} \\
Deskripsi Singkat & Kecerdasan Artifisial membangun pemahaman dan kemampuan mengimplementasikan algoritma pencarian, representasi pengetahuan, sistem pakar, dan pengantar ML untuk menyelesaikan permasalahan komputasi cerdas. \\
Bahan Kajian / Materi Pembelajaran & \begin{enumerate}\item Pengantar AI, agen cerdas, dan algoritma pencarian\item Representasi pengetahuan, logika, dan sistem inferensi\item Sistem pakar, Jaringan Bayes, dan probabilistic reasoning\item Machine learning, NLP dasar, dan Computer Vision dasar\end{enumerate} \\
Pustaka Utama & \begin{enumerate}\item Russell, S. \& Norvig, P. 2021. Artificial Intelligence: A Modern Approach (4th ed.). Pearson.\item Ertel, W. 2017. Introduction to Artificial Intelligence. Springer.\item Nilsson, N. 2010. The Quest for Artificial Intelligence. Cambridge University Press.\end{enumerate} \\
Pustaka Pendukung &  Materi LMS, dokumentasi resmi, dan jobsheet praktikum. \\
\end{longtblr}

\subsubsection*{Jaringan Komputer (RTI254005)}
\begin{longtblr}{width=\textwidth,colspec={Q[l,m,3cm]X[l,m]},rowsep=2pt,hlines,vlines,hspan=minimal}
SKS / Jam & 2 SKS/4 jam \\
Semester & 4 \\
CPL & \begin{itemize}\item \textbf{CPL05}: Mampu menghasilkan solusi teknologi komputasi dan infrastruktur digital yang sesuai dengan kebutuhan organisasi dan pengguna.\item \textbf{CPL09}: Mampu menjelaskan cara kerja sistem komputer dan menerapkan pendekatan komputasi untuk menyelesaikan masalah.\item \textbf{CPL10}: Mampu menganalisis persoalan kompleks dengan wawasan multidisiplin untuk menghasilkan solusi di bidang informatika dan ilmu komputer.\end{itemize} \\
Deskripsi Singkat & Jaringan Komputer membangun pemahaman komprehensif tentang arsitektur jaringan, protokol komunikasi, perangkat jaringan, dan keamanan jaringan untuk merancang infrastruktur jaringan yang efektif dan aman. \\
Bahan Kajian / Materi Pembelajaran & \begin{enumerate}\item Pengantar jaringan, model OSI, TCP/IP, dan media transmisi\item Protokol application, transport, dan network layer\item Wireless networking dan data link layer\item Keamanan jaringan: ancaman, firewall, VPN, dan monitoring\end{enumerate} \\
Pustaka Utama & \begin{enumerate}\item Tanenbaum, A. 2021. Computer Networks (6th ed.). Pearson.\item Forouzan, B. 2018. Data Communications and Networking. McGraw-Hill.\item Kurose, J. \& Ross, K. 2021. Computer Networking: A Top-Down Approach. Pearson.\end{enumerate} \\
Pustaka Pendukung &  Materi LMS, dokumentasi resmi, dan jobsheet praktikum. \\
\end{longtblr}

\subsubsection*{Praktikum Jaringan Komputer (RTI254006)}
\begin{longtblr}{width=\textwidth,colspec={Q[l,m,3cm]X[l,m]},rowsep=2pt,hlines,vlines,hspan=minimal}
SKS / Jam & 2 SKS/4 jam \\
Semester & 4 \\
CPL & \begin{itemize}\item \textbf{CPL05}: Mampu menyediakan solusi teknologi komputasi dan infrastruktur digital yang memenuhi kebutuhan pengguna dan organisasi.\item \textbf{CPL09}: Mampu memahami cara kerja sistem komputer dan menerapkan pendekatan komputasi untuk problem solving.\end{itemize} \\
Deskripsi Singkat & Praktikum Jaringan Komputer memberikan pengalaman langsung dalam mengkonfigurasi media jaringan, protokol TCP/IP, subnetting IPv4, routing, dan simulasi jaringan menggunakan Cisco Packet Tracer. \\
Bahan Kajian / Materi Pembelajaran & \begin{enumerate}\item Media dan konfigurasi perangkat jaringan\item Protokol lapisan aplikasi dan transport\item IPv4, subnetting, dan routing\item Packet Tracer dan analisis traffic\end{enumerate} \\
Pustaka Utama & \begin{enumerate}\item Cisco Networking Academy. 2020. CCNA: Introduction to Networks. Cisco Press.\item Forouzan, B.A. 2012. Data Communications and Networking. McGraw-Hill.\item Modul praktikum jaringan komputer JTI Polinema.\end{enumerate} \\
Pustaka Pendukung &  Materi LMS, dokumentasi resmi perangkat lunak, dan jobsheet praktikum. \\
\end{longtblr}

\subsubsection*{Pemrograman Web Lanjut (RTI254007)}
\begin{longtblr}{width=\textwidth,colspec={Q[l,m,3cm]X[l,m]},rowsep=2pt,hlines,vlines,hspan=minimal}
SKS / Jam & 3 SKS/6 jam \\
Semester & 4 \\
CPL & \begin{itemize}\item \textbf{CPL02}: Mampu menerapkan prinsip rekayasa sistem dan perangkat lunak untuk menghasilkan solusi aplikasi multi-platform sesuai kebutuhan.\item \textbf{CPL06}: Mampu merancang, mengimplementasikan, menguji, melakukan deployment, memelihara, serta menjamin mutu perangkat lunak yang menjawab permasalahan dan memenuhi kebutuhan pengguna.\item \textbf{CPL10}: Mampu menganalisis persoalan kompleks dengan wawasan multidisiplin untuk merumuskan solusi informatika atau ilmu komputer.\end{itemize} \\
Deskripsi Singkat & Mata kuliah Pemrograman Web Lanjut membahas pengembangan aplikasi web modern menggunakan Laravel, mencakup MVC, routing, migration, ORM, Blade, validasi, authentication, authorization, import/export, API, testing, dan evaluasi kualitas aplikasi. \\
Bahan Kajian / Materi Pembelajaran & \begin{enumerate}\item Laravel, MVC, routing, controller, view, dan Git\item Migration, Eloquent ORM, Blade, layout, dan validasi\item Authentication, authorization, import/export, API, testing, dan proyek Laravel\end{enumerate} \\
Pustaka Utama & \begin{enumerate}\item Azamuddin, M. and Mukhlasin, H. 2019. Laravel the PHP Framework for Web Artisans.\item Laravel Documentation. 2024. https://laravel.com/docs/10.x.\end{enumerate} \\
Pustaka Pendukung &  Materi LMS, dokumentasi resmi perangkat lunak, dan jobsheet praktikum. \\
\end{longtblr}

\subsubsection*{Statistik Komputasi (RTI254008)}
\begin{longtblr}{width=\textwidth,colspec={Q[l,m,3cm]X[l,m]},rowsep=2pt,hlines,vlines,hspan=minimal}
SKS / Jam & 2 SKS/4 jam \\
Semester & 4 \\
CPL & \begin{itemize}\item \textbf{CPL07}: Mampu mengembangkan sistem komputasi cerdas dan melakukan analisis data berbasis kecerdasan buatan untuk pengambilan keputusan.\item \textbf{CPL09}: Mampu menjelaskan cara kerja sistem komputer dan menerapkan pendekatan komputasi untuk menyelesaikan masalah.\item \textbf{CPL10}: Mampu menganalisis persoalan kompleks dengan wawasan multidisiplin untuk menghasilkan solusi di bidang informatika dan ilmu komputer.\end{itemize} \\
Deskripsi Singkat & Statistik Komputasi membangun kemampuan mengaplikasikan metode statistika deskriptif, inferensial, regresi, dan simulasi menggunakan Python/R untuk analisis data dan pengambilan keputusan berbasis bukti. \\
Bahan Kajian / Materi Pembelajaran & \begin{enumerate}\item Statistika deskriptif dan probabilitas\item Distribusi dan inferensi statistik\item Regresi, ANOVA, dan uji non-parametrik\item Simulasi Monte Carlo dan analisis data eksploratori\item Visualisasi data dan time series dasar\end{enumerate} \\
Pustaka Utama & \begin{enumerate}\item Walpole, R. E. 2016. Probability \& Statistics for Engineers and Scientists (9th ed.). Pearson.\item VanderPlas, J. 2016. Python Data Science Handbook. O'Reilly.\item Montgomery, D. 2017. Applied Statistics and Probability for Engineers. Wiley.\end{enumerate} \\
Pustaka Pendukung &  Materi LMS, dokumentasi resmi Python/R, dan jobsheet praktikum. \\
\end{longtblr}

\subsubsection*{Proyek Sistem Informasi (RTI254009)}
\begin{longtblr}{width=\textwidth,colspec={Q[l,m,3cm]X[l,m]},rowsep=2pt,hlines,vlines,hspan=minimal}
SKS / Jam & 3 SKS/6 jam \\
Semester & 4 \\
CPL & \begin{itemize}\item \textbf{CPL02}: Mampu menerapkan prinsip rekayasa sistem dan perangkat lunak untuk menghasilkan solusi aplikasi multi-platform sesuai kebutuhan.\item \textbf{CPL06}: Mampu merancang, mengimplementasikan, menguji, melakukan deployment, memelihara, serta menjamin mutu perangkat lunak yang menjawab permasalahan dan memenuhi kebutuhan pengguna.\item \textbf{CPL08}: Mampu mengelola proyek teknologi informasi secara profesional melalui perencanaan, koordinasi, komunikasi, dan optimalisasi sumber daya.\end{itemize} \\
Deskripsi Singkat & Proyek Sistem Informasi adalah mata kuliah berbasis proyek yang mengintegrasikan kemampuan analisis, desain, implementasi, dan manajemen proyek dalam pengembangan SI berbasis web secara tim menggunakan metodologi agile. \\
Bahan Kajian / Materi Pembelajaran & \begin{enumerate}\item Perencanaan proyek: WBS, jadwal, anggaran, risiko\item Metodologi SDLC dan Scrum/Agile\item Implementasi SI berbasis web: back-end, front-end, dan database\item Sprint execution, review, dan retrospective\item Deployment, UAT, laporan akhir, dan demo proyek\end{enumerate} \\
Pustaka Utama & \begin{enumerate}\item Sommerville, I. 2016. Software Engineering (10th ed.). Pearson.\item Schwaber, K. 2020. The Scrum Guide. Scrum.org.\item PMI. 2021. PMBOK Guide (7th ed.).\end{enumerate} \\
Pustaka Pendukung &  Materi LMS, dokumentasi resmi framework web, dan jobsheet praktikum. \\
\end{longtblr}

\subsection*{Semester 5}

\subsubsection*{Kewirausahaan Berbasis Teknologi (RTI255001)}
\begin{longtblr}{width=\textwidth,colspec={Q[l,m,3cm]X[l,m]},rowsep=2pt,hlines,vlines,hspan=minimal}
SKS / Jam & 2 SKS/4 jam \\
Semester & 5 \\
CPL & \begin{itemize}\item \textbf{CPL01}: Mampu menerapkan nilai ketakwaan, kepatuhan hukum, disiplin, profesionalisme, etika profesi, dan kewirausahaan dalam praktik profesi teknologi informasi.\item \textbf{CPL03}: Mampu mengelola diri, bekerja profesional, berkomunikasi lisan dan tulisan, serta melakukan presentasi secara efektif.\end{itemize} \\
Deskripsi Singkat & Mata kuliah Kewirausahaan Berbasis Teknologi membangun jiwa dan kompetensi kewirausahaan mahasiswa melalui pemahaman konsep kewirausahaan, simulasi bisnis (Business Game), analisis pasar, perencanaan keuangan, penyusunan Business Model Canvas, serta pengembangan dan presentasi rencana bisnis berbasis teknologi informasi. \\
Bahan Kajian / Materi Pembelajaran & \begin{enumerate}\item Konsep kewirausahaan, karakteristik wirausaha, dan intensi kewirausahaan\item Analisis pasar, aspek teknis bisnis, dan manajemen keuangan dasar\item Business Model Canvas, sumber pendanaan, dan aspek legal bisnis\item Manajemen SDM dan pengembangan produk baru\item Penyusunan dan presentasi rencana bisnis teknologi\end{enumerate} \\
Pustaka Utama & \begin{enumerate}\item Hisrich, R.D., Peters, M.P., Shepherd, D.A. 2017. Entrepreneurship. 10th ed. McGraw-Hill Education.\item Osterwalder, A., Pigneur, Y. 2010. Business Model Generation. John Wiley \& Sons.\item Bygrave, W., Zacharakis, A. 2011. Entrepreneurship. 2nd ed. John Wiley \& Sons.\end{enumerate} \\
Pustaka Pendukung &  Materi LMS, modul business game, jurnal kewirausahaan, dan referensi dosen. \\
\end{longtblr}

\subsubsection*{Metodologi Penelitian (RTI255002)}
\begin{longtblr}{width=\textwidth,colspec={Q[l,m,3cm]X[l,m]},rowsep=2pt,hlines,vlines,hspan=minimal}
SKS / Jam & 2 SKS/4 jam \\
Semester & 5 \\
CPL & \begin{itemize}\item \textbf{CPL04}: Mampu mendeskripsikan secara saintifik hasil kajian, pengembangan, dan implementasi TIK dalam bentuk skripsi, laporan, atau artikel ilmiah.\item \textbf{CPL03}: Mampu mengelola tim, manajemen diri, berkomunikasi lisan/tulisan, dan mempresentasikan dalam konteks profesional teknologi informasi.\end{itemize} \\
Deskripsi Singkat & Mata kuliah Metodologi Penelitian membangun kemampuan mahasiswa dalam merencanakan, merancang, dan melaksanakan penelitian TI secara ilmiah, mencakup studi literatur, identifikasi masalah, penyusunan metode, analisis data, dan penulisan proposal penelitian/skripsi. \\
Bahan Kajian / Materi Pembelajaran & \begin{enumerate}\item Jenis penelitian dan etika ilmiah\item Studi literatur dan identifikasi masalah\item Metode penelitian dan pengumpulan data\item Analisis data dan pemodelan sistem\item Penulisan proposal penelitian\end{enumerate} \\
Pustaka Utama & \begin{enumerate}\item Sugiyono. 2019. Metode Penelitian Kuantitatif, Kualitatif, dan R\&D. Alfabeta.\item Creswell, J.W. 2014. Research Design: Qualitative, Quantitative, and Mixed Methods. SAGE.\item Panduan Penulisan Laporan Akhir/Skripsi JTI Polinema.\end{enumerate} \\
Pustaka Pendukung &  Materi LMS, artikel ilmiah dari jurnal nasional/internasional, dan jobsheet praktikum. \\
\end{longtblr}

\subsubsection*{Pemrograman Mobile (RTI255003)}
\begin{longtblr}{width=\textwidth,colspec={Q[l,m,3cm]X[l,m]},rowsep=2pt,hlines,vlines,hspan=minimal}
SKS / Jam & 3 SKS/6 jam \\
Semester & 5 \\
CPL & \begin{itemize}\item \textbf{CPL02}: Mampu menerapkan prinsip rekayasa sistem dan perangkat lunak untuk menghasilkan solusi aplikasi multi-platform sesuai kebutuhan.\item \textbf{CPL06}: Mampu merancang, mengimplementasikan, menguji, melakukan deployment, memelihara, serta menjamin mutu perangkat lunak yang menjawab permasalahan dan memenuhi kebutuhan pengguna.\item \textbf{CPL10}: Mampu menganalisis persoalan kompleks dengan wawasan multidisiplin untuk merumuskan solusi informatika atau ilmu komputer.\end{itemize} \\
Deskripsi Singkat & Mata kuliah Pemrograman Mobile membahas pengembangan aplikasi mobile dengan Dart dan Flutter, mulai dari dasar bahasa, widget, layout, navigasi, plugin, kamera, state management, asynchronous programming, stream, persistence, RESTful API, hingga proyek aplikasi mobile. \\
Bahan Kajian / Materi Pembelajaran & \begin{enumerate}\item Dasar Dart dan Flutter\item Widget, layout, navigasi, plugin, dan kamera\item State management, asynchronous programming, stream, persistence, dan RESTful API\item Proyek aplikasi mobile berbasis Flutter\end{enumerate} \\
Pustaka Utama & \begin{enumerate}\item Flutter Documentation. https://docs.flutter.dev.\item Dart Documentation. https://dart.dev/guides.\item Windmill, E. Flutter in Action. Manning.\end{enumerate} \\
Pustaka Pendukung &  Materi LMS, dokumentasi resmi perangkat lunak, dan jobsheet praktikum. \\
\end{longtblr}

\subsubsection*{Pembelajaran Mesin (RTI255004)}
\begin{longtblr}{width=\textwidth,colspec={Q[l,m,3cm]X[l,m]},rowsep=2pt,hlines,vlines,hspan=minimal}
SKS / Jam & 3 SKS/6 jam \\
Semester & 5 \\
CPL & \begin{itemize}\item \textbf{CPL07}: Mampu mengembangkan sistem komputasi cerdas dan melakukan analisis data berbasis kecerdasan buatan untuk pengambilan keputusan.\item \textbf{CPL10}: Mampu menganalisis persoalan kompleks dengan wawasan multidisiplin untuk menghasilkan solusi di bidang informatika dan ilmu komputer.\end{itemize} \\
Deskripsi Singkat & Pembelajaran Mesin membangun kemampuan mengimplementasikan, mengevaluasi, dan mengoptimalkan algoritma ML untuk klasifikasi, regresi, clustering, dan deep learning dasar menggunakan scikit-learn dan Python dalam proyek end-to-end. \\
Bahan Kajian / Materi Pembelajaran & \begin{enumerate}\item Supervised learning: regresi, decision tree, random forest, SVM, k-NN\item Model evaluation: cross-validation, confusion matrix, metrik, regularisasi\item Unsupervised learning: k-Means, DBSCAN, dan dimensionality reduction (PCA)\item Ensemble methods, deep learning dasar, dan interpretabilitas model\item End-to-end ML project: pipeline, optimasi, dan presentasi\end{enumerate} \\
Pustaka Utama & \begin{enumerate}\item Géron, A. 2022. Hands-On Machine Learning (3rd ed.). O'Reilly.\item Raschka, S. 2022. Machine Learning with PyTorch and Scikit-Learn. Packt.\item James, G. 2021. An Introduction to Statistical Learning. Springer.\end{enumerate} \\
Pustaka Pendukung &  Materi LMS, dokumentasi resmi scikit-learn dan PyTorch, dan jobsheet praktikum. \\
\end{longtblr}

\subsubsection*{Business Intelligence (RTI255005)}
\begin{longtblr}{width=\textwidth,colspec={Q[l,m,3cm]X[l,m]},rowsep=2pt,hlines,vlines,hspan=minimal}
SKS / Jam & 2 SKS/4 jam \\
Semester & 5 \\
CPL & \begin{itemize}\item \textbf{CPL07}: Mampu mengembangkan sistem komputasi cerdas dan melakukan analisis data berbasis kecerdasan buatan untuk pengambilan keputusan.\item \textbf{CPL10}: Mampu menganalisis persoalan kompleks dengan wawasan multidisiplin untuk menghasilkan solusi di bidang informatika dan ilmu komputer.\end{itemize} \\
Deskripsi Singkat & Mata kuliah Business Intelligence membangun kemampuan merancang data warehouse, proses ETL, analisis OLAP, dan visualisasi dashboard interaktif untuk menghasilkan insight bisnis yang mendukung pengambilan keputusan strategis. \\
Bahan Kajian / Materi Pembelajaran & \begin{enumerate}\item Pengantar BI, data warehouse, dan proses ETL\item Analisis OLAP dan operasi multidimensional\item Visualisasi data, tools BI, dan desain dashboard\item Data mining, predictive analytics, dan real-time BI\item Implementasi proyek dashboard BI dan evaluasi insight\end{enumerate} \\
Pustaka Utama & \begin{enumerate}\item Kimball, R. \& Ross, M. 2013. The Data Warehouse Toolkit (3rd ed.). Wiley.\item Few, S. 2012. Show Me the Numbers. Analytics Press.\item Dokumentasi Power BI dan Tableau resmi.\end{enumerate} \\
Pustaka Pendukung &  Materi LMS, dokumentasi resmi perangkat lunak, dan jobsheet praktikum. \\
\end{longtblr}

\subsubsection*{Penjaminan Mutu Perangkat Lunak (RTI255006)}
\begin{longtblr}{width=\textwidth,colspec={Q[l,m,3cm]X[l,m]},rowsep=2pt,hlines,vlines,hspan=minimal}
SKS / Jam & 2 SKS/4 jam \\
Semester & 5 \\
CPL & \begin{itemize}\item \textbf{CPL06}: Mampu merancang, mengimplementasikan, menguji, melakukan deployment, memelihara, serta menjamin mutu perangkat lunak yang menjawab permasalahan dan memenuhi kebutuhan pengguna.\end{itemize} \\
Deskripsi Singkat & Mata kuliah Penjaminan Mutu Perangkat Lunak memberikan pengetahuan, keterampilan, dan pengalaman terarah untuk merancang strategi penjaminan mutu perangkat lunak melalui pengujian, verifikasi, validasi, dokumentasi pengujian, functional testing, non-functional testing, serta penyusunan rekomendasi peningkatan kualitas perangkat lunak sesuai standar industri dan kebutuhan pengguna. \\
Bahan Kajian / Materi Pembelajaran & \begin{enumerate}\item Pengantar SQA dan standar ISO.\item Pendekatan, tipe, dan level pengujian.\item Test planning, test scenario, test case, dan traceability.\item Test monitoring, test control, test completion, dan test reporting.\item Functional testing.\item Non-functional testing.\item Strategi QA komprehensif.\end{enumerate} \\
Pustaka Utama & \begin{enumerate}\item Sommerville, I. 2016. \textit{Software Engineering}, 10th Edition. Pearson.\item Myers, G. J., Sandler, C., and Badgett, T. 2011. \textit{The Art of Software Testing}, 3rd Edition. Wiley.\item Pressman, R. S. and Maxim, B. R. 2019. \textit{Software Engineering: A Practitioner's Approach}, 9th Edition. McGraw-Hill.\end{enumerate} \\
Pustaka Pendukung & \begin{enumerate}\item ISO/IEC/IEEE 29119 Software Testing Standard.\item ISO/IEC 25010 Systems and Software Quality Requirements and Evaluation.\item ISTQB Certified Tester Foundation Level Syllabus.\end{enumerate} \\
\end{longtblr}

\subsubsection*{Pengolahan Citra dan Visi Komputer (RTI255007)}
\begin{longtblr}{width=\textwidth,colspec={Q[l,m,3cm]X[l,m]},rowsep=2pt,hlines,vlines,hspan=minimal}
SKS / Jam & 3 SKS/6 jam \\
Semester & 5 \\
CPL & \begin{itemize}\item \textbf{CPL07}: Mampu mengembangkan sistem komputasi cerdas dan melakukan analisis data berbasis kecerdasan buatan untuk pengambilan keputusan.\item \textbf{CPL10}: Mampu menganalisis persoalan kompleks dengan wawasan multidisiplin untuk menghasilkan solusi di bidang informatika dan ilmu komputer.\end{itemize} \\
Deskripsi Singkat & Mata kuliah Pengolahan Citra dan Visi Komputer membangun kemampuan memproses citra digital dengan OpenCV dan mengembangkan sistem visi komputer berbasis CNN untuk klasifikasi, deteksi objek, dan segmentasi pada aplikasi nyata. \\
Bahan Kajian / Materi Pembelajaran & \begin{enumerate}\item Pengolahan citra digital: representasi, operasi pixel, histogram, dan color space\item Filtering spasial, deteksi tepi, dan segmentasi citra\item Transformasi morfologi dan evaluasi kualitas citra\item CNN: arsitektur, klasifikasi, deteksi objek, dan segmentasi gambar\item Transfer learning, augmentasi data, dan evaluasi model visi komputer\end{enumerate} \\
Pustaka Utama & \begin{enumerate}\item Gonzalez, R. 2018. Digital Image Processing (4th ed.). Pearson.\item Goodfellow, I. 2016. Deep Learning. MIT Press.\item Dokumentasi OpenCV dan PyTorch torchvision.\end{enumerate} \\
Pustaka Pendukung &  Materi LMS, dokumentasi resmi perangkat lunak, dan jobsheet praktikum. \\
\end{longtblr}

\subsubsection*{Administrasi dan Keamanan Jaringan (RTI255008)}
\begin{longtblr}{width=\textwidth,colspec={Q[l,m,3cm]X[l,m]},rowsep=2pt,hlines,vlines,hspan=minimal}
SKS / Jam & 3 SKS/6 jam \\
Semester & 5 \\
CPL & \begin{itemize}\item \textbf{CPL05}: Solusi teknologi komputasi dan infrastruktur digital sesuai kebutuhan.\end{itemize} \\
Deskripsi Singkat & Mata kuliah ini membekali mahasiswa dengan pengetahuan dan keterampilan komprehensif untuk menjadi administrator jaringan yang kompeten dalam mengelola keamanan jaringan komputer di organisasi. Materi meliputi desain sistem keamanan, instalasi dan konfigurasi gateway internet, administrasi sistem jaringan, optimasi kinerja jaringan, dan integrasi SIEM, sesuai Standar Kompetensi Kerja Nasional Indonesia (SKKNI) bidang Jaringan Komputer dan Sistem Administrasi. \\
Bahan Kajian / Materi Pembelajaran & \begin{enumerate}\item Pengenalan infrastruktur \& keamanan jaringan, CIA Triad, dan penilaian risiko\item Identifikasi ancaman siber dan kerentanan sistem jaringan\item Desain topologi jaringan aman: segmentasi, VLAN, DMZ\item Simulasi desain keamanan jaringan\item Konsep dan pemilihan gateway internet\item Instalasi dan konfigurasi dasar gateway internet\item Konfigurasi NAT dan firewall pada gateway\item Administrasi sistem jaringan: AAA, RADIUS/TACACS+\item Manajemen pengguna, RBAC, dan kontrol akses\item Monitoring jaringan dan dasar logging/SIEM\item Simulasi administrasi dan troubleshooting jaringan\item Identifikasi bottleneck dan optimasi kinerja jaringan\item Benchmarking, QoS, dan traffic shaping\item Simulasi optimasi kinerja dan integrasi SIEM\end{enumerate} \\
Pustaka Utama & \begin{enumerate}\item Standar Kompetensi Kerja Nasional Indonesia (SKKNI) Sektor Komunikasi dan Informasi, Bidang Jaringan Komputer dan Sistem Administrasi.\item Stallings, W. Cryptography and Network Security: Principles and Practice. Pearson.\item Cisco Networking Academy. CCNA Security Course Booklet. Cisco Press.\item NIST Special Publications (SP 800 Series) terkait Keamanan Jaringan.\end{enumerate} \\
Pustaka Pendukung &  Materi ajar/modul dari dosen/tim pengajar; dokumentasi resmi GNS3 dan VirtualBox. \\
\end{longtblr}

\subsection*{Semester 6}

\subsubsection*{Komunikasi dan Etika Profesi (RTI256001)}
\begin{longtblr}{width=\textwidth,colspec={Q[l,m,3cm]X[l,m]},rowsep=2pt,hlines,vlines,hspan=minimal}
SKS / Jam & 2 SKS/4 jam \\
Semester & 6 \\
CPL & \begin{itemize}\item \textbf{CPL01}: Mampu menerapkan nilai ketakwaan, kesadaran hukum, kedisiplinan, profesionalisme, etika profesi, dan pembelajaran sepanjang hayat, serta tanggap terhadap isu sosial dan perkembangan teknologi.\item \textbf{CPL03}: Mampu mengelola tim dan diri sendiri secara profesional serta mengomunikasikan dan mempresentasikan gagasan maupun hasil kerja secara efektif baik lisan maupun tulisan.\end{itemize} \\
Deskripsi Singkat & Mata kuliah Komunikasi dan Etika Profesi membangun kompetensi etika profesi TI, komunikasi profesional, dan perencanaan karir melalui studi kasus, simulasi wawancara, dan penyusunan portofolio digital yang siap kerja. \\
Bahan Kajian / Materi Pembelajaran & \begin{enumerate}\item Etika profesi TI, kode etik ACM/IEEE, hukum ITE, dan K3\item Komunikasi profesional: presentasi teknis dan penulisan laporan\item Dinamika tim, kolaborasi lintas fungsi, dan isu sosial teknologi\item Personal branding, portofolio digital, dan perencanaan karir\item Strategi pencarian kerja, CV, wawancara, dan etika penelitian\end{enumerate} \\
Pustaka Utama & \begin{enumerate}\item ACM/IEEE Code of Ethics 2018.\item Locker, K. 2020. Business Communication. McGraw-Hill.\item Floridi, L. 2014. The Ethics of AI. MIT Press.\end{enumerate} \\
Pustaka Pendukung &  Materi LMS, dokumentasi resmi, dan jobsheet praktikum. \\
\end{longtblr}

\subsubsection*{Pengembangan Karir (RTI256002)}
\begin{longtblr}{width=\textwidth,colspec={Q[l,m,3cm]X[l,m]},rowsep=2pt,hlines,vlines,hspan=minimal}
SKS / Jam & 2 SKS/4 jam \\
Semester & 6 \\
CPL & \begin{itemize}\item \textbf{CPL01}: Mampu menunjukkan nilai ketakwaan, kedisiplinan, profesionalisme, etika profesi, dan komitmen terhadap pembelajaran berkelanjutan dalam konteks sosial dan teknologi.\item \textbf{CPL03}: Mampu mengelola diri, bekerja sama dalam tim, berkomunikasi secara lisan dan tulisan, serta mempresentasikan gagasan secara profesional dalam lingkungan kerja.\end{itemize} \\
Deskripsi Singkat & Mata kuliah ini bertujuan agar mahasiswa memahami dan mampu mengelola pengembangan karir secara sistematis sehingga dapat mencapai kepuasan dan kesuksesan dalam berkarir. Materi meliputi pemetaan minat dan potensi diri, profesi bidang informatika, personal branding, komunikasi profesional, manajemen dan perencanaan karir, serta penyusunan Career Development Plan (CDP) melalui pendekatan individual project dan presentasi profesional. \\
Bahan Kajian / Materi Pembelajaran & \begin{enumerate}\item Pengantar pengembangan karir: konsep, tahapan, dan faktor karir\item Pemetaan minat, bakat, dan potensi diri (self-assessment)\item Karir mandiri dan fleksibel di era digital (gig economy, freelance, startup)\item Profesi bidang informatika: klasifikasi, kompetensi, dan jalur karir\item Karir organisasi vs freelance: perbandingan dan strategi\item Karir era Industri 4.0: otomasi, AI, IoT, dan lifelong learning\item Personal branding: identitas digital, LinkedIn, dan reputasi profesional\item Public speaking profesional: komunikasi verbal/nonverbal dan presentasi karir\item Manajemen karir jangka panjang dan networking profesional\item Perencanaan karir: SMART goals, roadmap, dan strategi pencapaian\item Penilaian kinerja: KPI, feedback, dan pengaruh terhadap karir\item Internasionalisasi karir: peluang kerja global dan persyaratan internasional\item Budaya manajerial global dan kolaborasi lintas budaya\item Pembuatan CV profesional yang relevan dan ATS-friendly\end{enumerate} \\
Pustaka Utama & \begin{enumerate}\item Widyanti, R. (2021). Manajemen Karir (Teori, Konsep dan Praktik). Media Sains Indonesia.\end{enumerate} \\
Pustaka Pendukung & \begin{enumerate}\item Sinambela, L. P. (2021). Manajemen Sumber Daya Manusia: Membangun tim kerja yang solid untuk meningkatkan kinerja. Bumi Aksara.\item Materi ajar/modul dari dosen/tim pengajar dan LMS.\end{enumerate} \\
\end{longtblr}

\subsubsection*{Pemrograman Berbasis Framework (RTI256003)}
\begin{longtblr}{width=\textwidth,colspec={Q[l,m,3cm]X[l,m]},rowsep=2pt,hlines,vlines,hspan=minimal}
SKS / Jam & 2 SKS/4 jam \\
Semester & 6 \\
CPL & \begin{itemize}\item \textbf{CPL02}: Mampu menerapkan prinsip rekayasa sistem dan perangkat lunak untuk menghasilkan solusi aplikasi multi-platform sesuai kebutuhan.\item \textbf{CPL06}: Mampu merancang, mengimplementasikan, menguji, melakukan deployment, memelihara, serta menjamin mutu perangkat lunak yang menjawab permasalahan dan memenuhi kebutuhan pengguna.\end{itemize} \\
Deskripsi Singkat & Mata kuliah Pemrograman Berbasis Framework membahas pengembangan aplikasi web modern berbasis framework JavaScript/TypeScript, mencakup setup, routing, styling, API routes, CSR, SSR, SSG, ISR, middleware, authentication, multi-role, optimasi, unit testing, dan proyek PBL. \\
Bahan Kajian / Materi Pembelajaran & \begin{enumerate}\item Setup, routing, styling, custom document\item API routes, CSR, SSR, SSG, ISR, middleware, route protection\item Authentication, multi-role, optimasi, unit testing, dan proyek PBL\end{enumerate} \\
Pustaka Utama & \begin{enumerate}\item Next.js Documentation. https://nextjs.org/docs.\item React Documentation. https://react.dev.\item Testing Library Documentation. https://testing-library.com.\end{enumerate} \\
Pustaka Pendukung &  Materi LMS, dokumentasi resmi perangkat lunak, dan jobsheet praktikum. \\
\end{longtblr}

\subsubsection*{Proyek Teknologi Terintegrasi (RTI256104)}
\begin{longtblr}{width=\textwidth,colspec={Q[l,m,3cm]X[l,m]},rowsep=2pt,hlines,vlines,hspan=minimal}
SKS / Jam & 6 SKS/12 jam \\
Semester & 6 \\
CPL & \begin{itemize}\item \textbf{CPL02}: Mampu menerapkan prinsip rekayasa sistem dan perangkat lunak untuk menghasilkan solusi aplikasi multi-platform sesuai kebutuhan.\item \textbf{CPL06}: Mampu merancang, mengimplementasikan, menguji, melakukan deployment, memelihara, serta menjamin mutu perangkat lunak yang menjawab permasalahan dan memenuhi kebutuhan pengguna.\item \textbf{CPL08}: Mampu mengelola proyek teknologi informasi secara profesional melalui perencanaan, koordinasi, komunikasi, dan optimalisasi sumber daya.\item \textbf{CPL10}: Mampu menganalisis persoalan kompleks dengan wawasan multidisiplin untuk menghasilkan solusi di bidang informatika dan ilmu komputer.\end{itemize} \\
Deskripsi Singkat & Proyek Teknologi Terintegrasi adalah mata kuliah capstone jalur reguler yang mengintegrasikan seluruh kompetensi rekayasa perangkat lunak dan manajemen proyek dalam pengembangan produk nyata multi-platform menggunakan metodologi agile selama satu semester. \\
Bahan Kajian / Materi Pembelajaran & \begin{enumerate}\item Manajemen proyek TI: perencanaan, WBS, jadwal, anggaran, dan risiko\item Rekayasa perangkat lunak agile: Scrum, sprint planning, backlog, dan velocity\item Pengembangan aplikasi multi-platform: integrasi front-end, back-end, dan basis data\item Monitoring proyek, kontrol kualitas, deployment, dan defense produk akhir\end{enumerate} \\
Pustaka Utama & \begin{enumerate}\item Schwaber, K. 2020. The Scrum Guide. Scrum.org.\item Sommerville, I. 2016. Software Engineering (10th ed.). Pearson.\item PMI. 2021. PMBOK Guide (7th ed.).\end{enumerate} \\
Pustaka Pendukung &  Materi LMS, repositori proyek, dan jobsheet capstone. \\
\end{longtblr}

\subsubsection*{Internet of Things (RTI256105)}
\begin{longtblr}{width=\textwidth,colspec={Q[l,m,3cm]X[l,m]},rowsep=2pt,hlines,vlines,hspan=minimal}
SKS / Jam & 3 SKS/6 jam \\
Semester & 6 \\
CPL & \begin{itemize}\item \textbf{CPL05}: Mampu menghasilkan solusi teknologi komputasi dan infrastruktur digital yang sesuai dengan kebutuhan organisasi dan pengguna.\item \textbf{CPL02}: Mampu menerapkan prinsip rekayasa sistem dan perangkat lunak untuk menghasilkan solusi aplikasi multi-platform sesuai kebutuhan.\item \textbf{CPL09}: Mampu menjelaskan cara kerja sistem komputer dan menerapkan pendekatan komputasi untuk menyelesaikan masalah.\end{itemize} \\
Deskripsi Singkat & Internet of Things membangun kompetensi merancang, mengimplementasikan, dan mendeploy sistem IoT end-to-end dari sensor hingga cloud menggunakan protokol standar industri dan platform IoT modern. \\
Bahan Kajian / Materi Pembelajaran & \begin{enumerate}\item Arsitektur IoT: konsep, komponen, sensor, aktuator, gateway, dan layer edge-fog-cloud\item Mikrokontroler dan SBC: Arduino, Raspberry Pi, dan antarmuka sensor\item Protokol komunikasi IoT: MQTT, HTTP, CoAP, WebSocket, dan pemrosesan data sensor\item Platform IoT: Node-RED, AWS IoT, keamanan, integrasi cloud, dan visualisasi data real-time\end{enumerate} \\
Pustaka Utama & \begin{enumerate}\item Buyya, R. \& Dastjerdi, A. 2016. Internet of Things: Principles and Paradigms. Elsevier.\item McEwen, A. 2014. Designing the Internet of Things. Wiley.\item Dokumentasi AWS IoT Core dan Node-RED.\end{enumerate} \\
Pustaka Pendukung &  Materi LMS, dokumentasi resmi platform IoT, dan jobsheet praktikum. \\
\end{longtblr}

\subsubsection*{Big Data (RTI256106)}
\begin{longtblr}{width=\textwidth,colspec={Q[l,m,3cm]X[l,m]},rowsep=2pt,hlines,vlines,hspan=minimal}
SKS / Jam & 2 SKS/4 jam \\
Semester & 6 \\
CPL & \begin{itemize}\item \textbf{CPL07}: Mampu mengembangkan sistem komputasi cerdas dan melakukan analisis data berbasis kecerdasan buatan untuk pengambilan keputusan.\item \textbf{CPL05}: Mampu menghasilkan solusi teknologi komputasi dan infrastruktur digital yang sesuai dengan kebutuhan organisasi dan pengguna.\item \textbf{CPL09}: Mampu menjelaskan cara kerja sistem komputer dan menerapkan pendekatan komputasi untuk menyelesaikan masalah.\end{itemize} \\
Deskripsi Singkat & Big Data membangun kemampuan merancang dan mengimplementasikan pipeline pemrosesan data skala besar menggunakan ekosistem Hadoop dan Spark, mencakup batch processing, stream processing, dan visualisasi data. \\
Bahan Kajian / Materi Pembelajaran & \begin{enumerate}\item Konsep Big Data: 4V (Volume, Velocity, Variety, Veracity), arsitektur Lambda, Kappa, dan HDFS\item Pemrosesan terdistribusi: MapReduce, Hadoop ekosistem (YARN, Hive, Pig), dan Apache Spark (RDD, DataFrame, SQL)\item Stream processing: Spark Streaming, Apache Kafka, dan pemrosesan event real-time\item Data lakehouse: Delta Lake, Apache Iceberg, NoSQL (Cassandra, MongoDB), dan visualisasi Big Data\end{enumerate} \\
Pustaka Utama & \begin{enumerate}\item White, T. 2015. Hadoop: The Definitive Guide. O'Reilly.\item Chambers, B. 2018. Spark: The Definitive Guide. O'Reilly.\item Marz, N. 2015. Big Data: Principles and Best Practices. Manning.\end{enumerate} \\
Pustaka Pendukung &  Materi LMS, dokumentasi resmi Apache Hadoop dan Spark, dan jobsheet praktikum. \\
\end{longtblr}

\subsubsection*{Cloud Computing (RTI256107)}
\begin{longtblr}{width=\textwidth,colspec={Q[l,m,3cm]X[l,m]},rowsep=2pt,hlines,vlines,hspan=minimal}
SKS / Jam & 2 SKS/4 jam \\
Semester & 6 \\
CPL & \begin{itemize}\item \textbf{CPL05}: Mampu menghasilkan solusi teknologi komputasi dan infrastruktur digital yang sesuai dengan kebutuhan organisasi dan pengguna.\item \textbf{CPL09}: Mampu menjelaskan cara kerja sistem komputer dan menerapkan pendekatan komputasi untuk menyelesaikan masalah.\item \textbf{CPL10}: Mampu menganalisis persoalan kompleks dengan wawasan multidisiplin untuk menghasilkan solusi di bidang informatika dan ilmu komputer.\end{itemize} \\
Deskripsi Singkat & Cloud Computing membangun kemampuan merancang, mengimplementasikan, dan mengevaluasi solusi cloud menggunakan IaaS/PaaS/SaaS, containerisasi Docker-Kubernetes, dan prinsip arsitektur cloud yang andal dan scalable. \\
Bahan Kajian / Materi Pembelajaran & \begin{enumerate}\item Virtualisasi, containerisasi, dan model layanan cloud (IaaS/PaaS/SaaS)\item Docker dan Kubernetes: containerisasi dan orkestrasi\item Arsitektur cloud: microservices, VPC, keamanan, dan jaringan\item Optimasi biaya, governance, dan strategi multi-cloud/hybrid\end{enumerate} \\
Pustaka Utama & \begin{enumerate}\item Velte, T. 2010. Cloud Computing: A Practical Approach. McGraw-Hill.\item Burns, B. 2019. Kubernetes: Up and Running. O'Reilly.\item Dokumentasi AWS/GCP/Azure resmi.\end{enumerate} \\
Pustaka Pendukung &  Materi LMS, dokumentasi resmi platform cloud, dan jobsheet praktikum. \\
\end{longtblr}

\subsubsection*{Komputasi Hijau (RTI256108)}
\begin{longtblr}{width=\textwidth,colspec={Q[l,m,3cm]X[l,m]},rowsep=2pt,hlines,vlines,hspan=minimal}
SKS / Jam & 2 SKS/4 jam \\
Semester & 6 \\
CPL & \begin{itemize}\item \textbf{CPL01}: Mampu menerapkan nilai ketakwaan, kesadaran hukum, kedisiplinan, profesionalisme, etika profesi, dan pembelajaran sepanjang hayat, serta tanggap terhadap isu sosial dan perkembangan teknologi.\item \textbf{CPL05}: Mampu menghasilkan solusi teknologi komputasi dan infrastruktur digital yang sesuai dengan kebutuhan organisasi dan pengguna.\end{itemize} \\
Deskripsi Singkat & Komputasi Hijau membangun kesadaran dan kemampuan merancang infrastruktur TI ramah lingkungan melalui efisiensi energi, virtualisasi, pengelolaan e-waste, dan penerapan prinsip keberlanjutan dalam pengembangan sistem TI. \\
Bahan Kajian / Materi Pembelajaran & \begin{enumerate}\item Konsep komputasi hijau, carbon footprint, dan regulasi keberlanjutan TI\item Efisiensi energi perangkat keras, PUE, dan data center hijau\item Virtualisasi, green software, cloud computing, dan renewable energy\item E-waste, green IT governance, dan studi kasus sustainabilitas\end{enumerate} \\
Pustaka Utama & \begin{enumerate}\item Murugesan, S. 2008. Harnessing Green IT. IEEE IT Professional.\item Belady, C. 2010. Green Grid Metrics. The Green Grid.\item Gupta, M. 2011. Green IT Strategy and Implementation. Springer.\end{enumerate} \\
Pustaka Pendukung &  Materi LMS, dokumentasi standar ISO 14001, laporan The Green Grid, dan jobsheet praktikum. \\
\end{longtblr}

\subsubsection*{Proyek Inovasi (RTI256204)}
\begin{longtblr}{width=\textwidth,colspec={Q[l,m,3cm]X[l,m]},rowsep=2pt,hlines,vlines,hspan=minimal}
SKS / Jam & 6 SKS/12 jam \\
Semester & 6 \\
CPL & \begin{itemize}\item \textbf{CPL08}: Mampu mengelola proyek teknologi informasi secara profesional melalui perencanaan, koordinasi, komunikasi, dan optimalisasi sumber daya.\item \textbf{CPL02}: Mampu menerapkan prinsip rekayasa sistem dan perangkat lunak untuk menghasilkan solusi aplikasi multi-platform sesuai kebutuhan.\item \textbf{CPL06}: Mampu merancang, mengimplementasikan, menguji, melakukan deployment, memelihara, serta menjamin mutu perangkat lunak yang menjawab permasalahan dan memenuhi kebutuhan pengguna.\end{itemize} \\
Deskripsi Singkat & Proyek Inovasi adalah capstone jalur magang/industri yang mengintegrasikan kreativitas, rekayasa, dan manajemen proyek dalam mengembangkan solusi TI inovatif berbasis masalah nyata industri menggunakan metodologi agile. \\
Bahan Kajian / Materi Pembelajaran & \begin{enumerate}\item Problem framing, opportunity analysis, dan technology landscape\item Proposal inovasi: WBS, jadwal, anggaran, dan manajemen risiko\item Metodologi agile/scrum: sprint planning, review, dan retrospective\item Implementasi MVP, UAT, deployment, dan dokumentasi proyek\end{enumerate} \\
Pustaka Utama & \begin{enumerate}\item Ries, E. 2011. The Lean Startup. Crown Business.\item Blank, S. 2013. The Startup Owner's Manual. K\&S Ranch.\item Schwaber, K. 2020. The Scrum Guide. Scrum.org.\end{enumerate} \\
Pustaka Pendukung &  Materi LMS, dokumentasi resmi perangkat lunak, dan jobsheet praktikum. \\
\end{longtblr}

\subsubsection*{Workshop Teknologi Terapan (RTI256205)}
\begin{longtblr}{width=\textwidth,colspec={Q[l,m,3cm]X[l,m]},rowsep=2pt,hlines,vlines,hspan=minimal}
SKS / Jam & 6 SKS/12 jam \\
Semester & 6 \\
CPL & \begin{itemize}\item \textbf{CPL02}: Mampu menerapkan prinsip rekayasa sistem dan perangkat lunak untuk menghasilkan solusi aplikasi multi-platform sesuai kebutuhan.\item \textbf{CPL06}: Mampu merancang, mengimplementasikan, menguji, melakukan deployment, memelihara, serta menjamin mutu perangkat lunak yang menjawab permasalahan dan memenuhi kebutuhan pengguna.\item \textbf{CPL08}: Mampu mengelola proyek teknologi informasi secara profesional melalui perencanaan, koordinasi, komunikasi, dan optimalisasi sumber daya.\end{itemize} \\
Deskripsi Singkat & Workshop Teknologi Terapan adalah capstone jalur magang/industri yang mengintegrasikan rekayasa perangkat lunak, quality assurance, dan manajemen proyek dalam mengembangkan produk teknologi nyata sesuai kebutuhan industri melalui pendekatan workshop berbasis sprint. \\
Bahan Kajian / Materi Pembelajaran & \begin{enumerate}\item Rekayasa perangkat lunak dalam konteks industri: sprint planning dan backlog management\item Pengembangan produk iteratif: implementasi fitur, integrasi, dan validasi\item Quality assurance: unit testing, integration testing, dan UAT\item Manajemen proyek workshop: stakeholder communication, demo, dan retrospective\end{enumerate} \\
Pustaka Utama & \begin{enumerate}\item Humble, J. 2011. Continuous Delivery. Pearson.\item Cohn, M. 2009. Succeeding with Agile. Addison-Wesley.\item Myers, G. 2011. The Art of Software Testing. Wiley.\end{enumerate} \\
Pustaka Pendukung &  Materi LMS, dokumentasi resmi perangkat lunak, dan jobsheet praktikum. \\
\end{longtblr}

\subsubsection*{Rekayasa Sistem (RTI256206)}
\begin{longtblr}{width=\textwidth,colspec={Q[l,m,3cm]X[l,m]},rowsep=2pt,hlines,vlines,hspan=minimal}
SKS / Jam & 6 SKS/12 jam \\
Semester & 6 \\
CPL & \begin{itemize}\item \textbf{CPL02}: Mampu menerapkan prinsip rekayasa sistem dan perangkat lunak untuk menghasilkan solusi aplikasi multi-platform sesuai kebutuhan.\item \textbf{CPL06}: Mampu merancang, mengimplementasikan, menguji, melakukan deployment, memelihara, serta menjamin mutu perangkat lunak yang menjawab permasalahan dan memenuhi kebutuhan pengguna.\item \textbf{CPL10}: Mampu menganalisis persoalan kompleks dengan wawasan multidisiplin untuk menghasilkan solusi di bidang informatika dan ilmu komputer.\end{itemize} \\
Deskripsi Singkat & Rekayasa Sistem adalah capstone jalur magang/industri yang mengintegrasikan rekayasa kebutuhan, perancangan arsitektur enterprise, dan rekayasa deployment untuk menghasilkan sistem yang handal, maintainable, dan siap produksi. \\
Bahan Kajian / Materi Pembelajaran & \begin{enumerate}\item Requirement elicitation: teknik wawancara, use case, user story, dan spesifikasi kebutuhan\item Arsitektur sistem enterprise: pola arsitektur, komponen, antarmuka, dan dokumentasi ADR\item Pipeline CI/CD: automated testing, integrasi, dan strategi deployment (blue-green, canary)\item Monitoring, observability, dan pemeliharaan sistem produksi\end{enumerate} \\
Pustaka Utama & \begin{enumerate}\item Sommerville, I. 2016. Software Engineering (10th ed.). Pearson.\item Bass, L. 2013. Software Architecture in Practice. Addison-Wesley.\item Forsgren, N. 2018. Accelerate. IT Revolution.\end{enumerate} \\
Pustaka Pendukung &  Materi LMS, dokumentasi resmi perangkat lunak, dan jobsheet praktikum. \\
\end{longtblr}

\subsubsection*{Teknologi Terapan (RTI256207)}
\begin{longtblr}{width=\textwidth,colspec={Q[l,m,3cm]X[l,m]},rowsep=2pt,hlines,vlines,hspan=minimal}
SKS / Jam & 6 SKS/12 jam \\
Semester & 6 \\
CPL & \begin{itemize}\item \textbf{CPL02}: Mampu menerapkan prinsip rekayasa sistem dan perangkat lunak untuk menghasilkan solusi aplikasi multi-platform sesuai kebutuhan.\item \textbf{CPL05}: Mampu menghasilkan solusi teknologi komputasi dan infrastruktur digital yang sesuai dengan kebutuhan organisasi dan pengguna.\item \textbf{CPL08}: Mampu mengelola proyek teknologi informasi secara profesional melalui perencanaan, koordinasi, komunikasi, dan optimalisasi sumber daya.\end{itemize} \\
Deskripsi Singkat & Teknologi Terapan adalah capstone jalur magang/industri yang mengintegrasikan analisis kebutuhan industri, perancangan, dan implementasi solusi teknologi tepat guna dalam konteks industri nyata melalui pendekatan iteratif dan kolaboratif bersama stakeholder. \\
Bahan Kajian / Materi Pembelajaran & \begin{enumerate}\item Needs assessment, technology mapping, dan feasibility study untuk solusi teknologi industri\item Perencanaan implementasi teknologi: roadmap, manajemen risiko, dan stakeholder management\item Implementasi solusi iteratif: prototipe, pengembangan, integrasi, dan optimasi\item Deployment produksi, transfer knowledge, dan evaluasi dampak teknologi\end{enumerate} \\
Pustaka Utama & \begin{enumerate}\item Rogers, E. 2003. Diffusion of Innovations. Free Press.\item Christensen, C. 2016. The Innovator's Dilemma. HBR Press.\item Brown, T. 2009. Change by Design. HarperBusiness.\end{enumerate} \\
Pustaka Pendukung &  Materi LMS, dokumentasi resmi perangkat lunak, dan jobsheet praktikum. \\
\end{longtblr}

\subsection*{Semester 7}

\subsubsection*{Magang (RTI257001)}
\begin{longtblr}{width=\textwidth,colspec={Q[l,m,3cm]X[l,m]},rowsep=2pt,hlines,vlines,hspan=minimal}
SKS / Jam & 20 SKS / 40 jam/minggu \\
Semester & 7 \\
CPL & \begin{itemize}\item \textbf{CPL03}: Mampu mengelola tim dan diri sendiri secara profesional serta mengomunikasikan dan mempresentasikan gagasan maupun hasil kerja secara efektif baik lisan maupun tulisan.\item \textbf{CPL06}: Mampu merancang, mengimplementasikan, menguji, melakukan deployment, memelihara, serta menjamin mutu perangkat lunak yang menjawab permasalahan dan memenuhi kebutuhan pengguna.\item \textbf{CPL08}: Mampu mengelola proyek teknologi informasi secara profesional melalui perencanaan, koordinasi, komunikasi, dan optimalisasi sumber daya.\end{itemize} \\
Deskripsi Singkat & Magang adalah program pembelajaran luar kampus (MBKM) 20 SKS yang menempatkan mahasiswa di industri TI selama ±4 bulan untuk mengaplikasikan kompetensi akademis, membangun profesionalisme, dan menghasilkan laporan ilmiah berbasis pengalaman kerja nyata. \\
Bahan Kajian / Materi Pembelajaran & \begin{enumerate}\item Pembekalan magang, etika profesi, dan orientasi industri\item Praktik kerja industri dan logbook kegiatan\item Aplikasi kompetensi rekayasa perangkat lunak di industri\item Penulisan laporan magang dan komunikasi profesional\item Presentasi dan sidang magang\end{enumerate} \\
Pustaka Utama & \begin{enumerate}\item Buku Pedoman Magang D4 TI Polinema 2025.\item Panduan MBKM Kemdikbud 2020.\item Pedoman Penulisan Laporan Magang Polinema.\end{enumerate} \\
Pustaka Pendukung &  Materi LMS, dokumen industri/pedoman magang, dan bahan bimbingan. \\
\end{longtblr}

\subsection*{Semester 8}

\subsubsection*{Skripsi (RTI258001)}
\begin{longtblr}{width=\textwidth,colspec={Q[l,m,3cm]X[l,m]},rowsep=2pt,hlines,vlines,hspan=minimal}
SKS / Jam & 8 SKS / 16 jam/minggu \\
Semester & 8 \\
CPL & \begin{itemize}\item \textbf{CPL04}: Mampu menyusun deskripsi saintifik hasil kajian, pengembangan, atau implementasi teknologi informasi dan komunikasi dalam bentuk skripsi, laporan, atau artikel ilmiah.\item \textbf{CPL06}: Mampu merancang, mengimplementasikan, menguji, melakukan deployment, memelihara, serta menjamin mutu perangkat lunak yang menjawab permasalahan dan memenuhi kebutuhan pengguna.\item \textbf{CPL08}: Mampu mengelola proyek teknologi informasi secara profesional melalui perencanaan, koordinasi, komunikasi, dan optimalisasi sumber daya.\item \textbf{CPL10}: Mampu menganalisis persoalan kompleks dengan wawasan multidisiplin untuk menghasilkan solusi di bidang informatika dan ilmu komputer.\end{itemize} \\
Deskripsi Singkat & Skripsi adalah mata kuliah tugas akhir 8 SKS yang mengintegrasikan seluruh kompetensi akademis mahasiswa dalam melaksanakan penelitian atau pengembangan TIK secara mandiri, menghasilkan karya ilmiah berstandar publikasi, dan mempertahankannya di hadapan tim penguji. \\
Bahan Kajian / Materi Pembelajaran & \begin{enumerate}\item Metodologi penelitian dan penyusunan proposal ilmiah\item Kajian literatur dan identifikasi masalah penelitian\item Pelaksanaan penelitian dan pengembangan sistem TIK\item Penulisan naskah skripsi sesuai standar ilmiah\item Seminar, sidang, dan diseminasi hasil penelitian\end{enumerate} \\
Pustaka Utama & \begin{enumerate}\item Buku Pedoman Skripsi D4 TI Polinema 2025.\item Creswell, J. 2018. Research Design. SAGE.\item Pressman, R. 2014. Software Engineering: A Practitioner's Approach. McGraw-Hill.\end{enumerate} \\
Pustaka Pendukung &  Materi LMS, jurnal ilmiah TIK, dan panduan penulisan karya ilmiah. \\
\end{longtblr}

\subsubsection*{Bahasa Inggris Persiapan Kerja (RTI258002)}
\begin{longtblr}{width=\textwidth,colspec={Q[l,m,3cm]X[l,m]},rowsep=2pt,hlines,vlines,hspan=minimal}
SKS / Jam & 2 SKS / 4 jam \\
Semester & 8 \\
CPL & \begin{itemize}\item \textbf{CPL01}: Mampu menerapkan nilai ketakwaan, kesadaran hukum, kedisiplinan, profesionalisme, etika profesi, dan pembelajaran sepanjang hayat, serta tanggap terhadap isu sosial dan perkembangan teknologi.\item \textbf{CPL03}: Mampu mengelola tim dan diri sendiri secara profesional serta mengomunikasikan dan mempresentasikan gagasan maupun hasil kerja secara efektif baik lisan maupun tulisan.\end{itemize} \\
Deskripsi Singkat & Bahasa Inggris Persiapan Kerja mempersiapkan mahasiswa memasuki pasar kerja TI global dengan membangun kemampuan komunikasi profesional dalam Bahasa Inggris mencakup penulisan dokumen karier, presentasi teknis, simulasi wawancara, dan komunikasi lintas budaya. \\
Bahan Kajian / Materi Pembelajaran & \begin{enumerate}\item Professional CV, LinkedIn profile, cover letter, dan email profesional\item Technical writing dan dokumentasi proyek TI\item Presentasi teknis dan komunikasi tim multikultural\item Job interview: STAR method, behavioral dan technical interview\item Business English, IT portfolio, dan mock interview\end{enumerate} \\
Pustaka Utama & \begin{enumerate}\item Mascull, B. 2017. \textit{Business Vocabulary in Use (Advanced)}. Cambridge.\item Guffey, M. 2019. \textit{Business Communication: Process and Product}. Cengage.\item Swan, M. 2016. \textit{Practical English Usage}. Oxford.\end{enumerate} \\
\end{longtblr}

\endgroup
